\documentclass[11pt,oneside]{article}

\usepackage[a4paper,margin=27mm,headheight=15pt]{geometry}
\usepackage[T1]{fontenc}
\usepackage[utf8]{inputenc}
\IfFileExists{libertinus.sty}{\usepackage{libertinus}}{\usepackage{lmodern}}
\usepackage{microtype}
\usepackage{amsmath,amssymb,amsthm,mathtools,mathrsfs}
% Canonical mathematical notation for active FabiusFunction LaTeX documents.
%
% This file is intentionally a plain TeX input rather than a LaTeX package.
% Every standalone document loads it by an explicit path relative to that
% document's directory; included fragments inherit the definitions from their
% root document.  Keep this file independent of document layout and metadata.
%
% Contract:
%   * one command has one mathematical meaning throughout the active corpus;
%   * names encode semantic roles rather than a document-local choice of letter;
%   * visually similar but differently normalized objects have different print;
%   * every command below is defined normatively in the notation catalogue.
%
% The active corpus excludes docs/archive/.  Historical sources may retain
% their original notation only inside an explicitly labelled quotation.

\makeatletter
\@ifundefined{mathbb}{%
  \PackageError{fabius-notation}{The amssymb package must be loaded first}%
    {Load amsmath and amssymb before inputting fabius-notation.tex.}%
}{}
\@ifundefined{operatorname}{%
  \PackageError{fabius-notation}{The amsmath package must be loaded first}%
    {Load amsmath before inputting fabius-notation.tex.}%
}{}
\makeatother

% Version stamp used by the catalogue and by static audits.
\newcommand{\FabiusNotationVersion}{2026-08-31}

% Number systems.  NaturalNumbers is explicitly the nonnegative convention.
\newcommand{\NaturalNumbers}{\mathbb{N}_{0}}
\newcommand{\PositiveIntegers}{\mathbb{N}_{>0}}
\newcommand{\IntegerNumbers}{\mathbb{Z}}
\newcommand{\RationalNumbers}{\mathbb{Q}}
\newcommand{\RealNumbers}{\mathbb{R}}
\newcommand{\ComplexNumbers}{\mathbb{C}}
\newcommand{\ExtendedRealNumbers}{\overline{\mathbb{R}}}

% The three principal Fabius--Rvachev functions.  Subscripts are deliberate:
% a bare F and a calligraphic F have several unrelated uses in the corpus.
\newcommand{\FabiusBounded}{F_{\mathrm{bd}}}
\newcommand{\FabiusGlobal}{F_{\mathrm{gl}}}
\newcommand{\FabiusQuantile}{Q_{F}}
\newcommand{\FabiusClampedQuantile}{Q_{F,\mathrm{cl}}}
\newcommand{\RvachevUp}{\operatorname{up}}

% Constants, elementary operators, and delimiters.
\newcommand{\EulerE}{\mathrm{e}}
\newcommand{\ImaginaryUnit}{\mathrm{i}}
\newcommand{\Differential}{\mathop{}\!\mathrm{d}}
\newcommand{\AbsoluteValue}[1]{\left\lvert #1\right\rvert}
\newcommand{\Norm}[1]{\left\lVert #1\right\rVert}
\newcommand{\Floor}[1]{\left\lfloor #1\right\rfloor}
\newcommand{\Ceiling}[1]{\left\lceil #1\right\rceil}
\newcommand{\FractionalPart}[1]{\left\{#1\right\}}
\newcommand{\PositivePart}[1]{\left(#1\right)_{+}}
\newcommand{\IndicatorOf}[1]{\mathbf{1}_{#1}}
\newcommand{\SetOf}[1]{\left\{#1\right\}}
\newcommand{\InnerProduct}[1]{\left\langle #1\right\rangle}
\newcommand{\CardinalityOf}[1]{\left\lvert #1\right\rvert}
\newcommand{\DefinitionEquals}{\mathrel{\coloneqq}}
\newcommand{\Sign}{\operatorname{sgn}}
\newcommand{\IdentityOperator}{\operatorname{id}}
\newcommand{\SpanOperator}{\operatorname{span}}
\newcommand{\TraceOperator}{\operatorname{tr}}
\newcommand{\RankOperator}{\operatorname{rank}}
\newcommand{\DiagonalOperator}{\operatorname{diag}}
\newcommand{\ResidueOperator}{\operatorname*{Res}}
\newcommand{\LeastCommonMultiple}{\operatorname{lcm}}
\newcommand{\OrderOperator}{\operatorname{ord}}
\newcommand{\PrincipalLogarithm}{\operatorname{Log}}
\newcommand{\FallingFactorial}[2]{\left(#1\right)^{\underline{#2}}}
\newcommand{\RisingFactorial}[2]{\left(#1\right)^{\overline{#2}}}

% Support, topology, convolution, and metric notation.
\newcommand{\SupportOperator}{\operatorname{supp}}
\newcommand{\SupportOf}[1]{\SupportOperator\!\left(#1\right)}
\newcommand{\TopologicalSupportOf}[1]{\operatorname{tsupp}\!\left(#1\right)}
\newcommand{\Convolution}{\mathbin{*}}
\newcommand{\MetricDistance}{\operatorname{dist}}
\newcommand{\TotalVariationDistance}{d_{\mathrm{TV}}}

% Probability.  Equality and convergence in law are not metric distance.
\newcommand{\Probability}{\mathbb{P}}
\newcommand{\Expectation}{\mathbb{E}}
\newcommand{\Variance}{\operatorname{Var}}
\newcommand{\Covariance}{\operatorname{Cov}}
\newcommand{\LawOperator}{\operatorname{Law}}
\newcommand{\LawOf}[1]{\LawOperator\!\left(#1\right)}
\newcommand{\EqualInLaw}{\mathrel{\overset{\mathrm{law}}{=}}}
\newcommand{\ConvergesInLaw}{\mathrel{\xrightarrow{\mathrm{law}}}}

% Fourier and integral transforms.  The printed labels expose normalization.
% FourierTwoPi uses exp(-2*pi*i*x*xi); FourierAngular uses exp(-i*x*omega).
\newcommand{\FourierTwoPi}[1]{\widehat{#1}_{\,2\pi}}
\newcommand{\FourierAngular}[1]{\widehat{#1}_{\,\mathrm{ang}}}
\newcommand{\LaplaceTransformSymbol}{\mathcal{L}}
\newcommand{\MellinTransformSymbol}{\mathcal{M}}
\newcommand{\LaplaceTransformOf}[1]{\LaplaceTransformSymbol\!\left[#1\right]}
\newcommand{\MellinTransformOf}[1]{\MellinTransformSymbol\!\left[#1\right]}
\newcommand{\MomentGeneratingFunction}[1]{M_{#1}}
\newcommand{\CharacteristicFunction}[1]{\varphi_{#1}}

% The two standard sinc functions must never share one symbol.
\newcommand{\SincRad}{\operatorname{sinc}_{\mathrm{rad}}}
\newcommand{\SincPi}{\operatorname{sinc}_{\pi}}

% Binary arithmetic and Thue--Morse notation.
\newcommand{\BinaryDigitSum}{\operatorname{wt}_{2}}
\newcommand{\ThueMorseSignSymbol}{\varepsilon}
\newcommand{\ThueMorseSign}[1]{\varepsilon_{#1}}
\newcommand{\PadicValuation}[1]{\nu_{#1}}
\newcommand{\TwoAdicValuation}{\nu_{2}}

% q-series notation.  Every base is an explicit argument.
\newcommand{\QPochhammer}[3]{\left(#1;#2\right)_{#3}}
\newcommand{\GaussianBinomial}[3]{%
  \genfrac{[}{]}{0pt}{}{#1}{#2}_{#3}%
}
\newcommand{\QInteger}[2]{\left[#1\right]_{#2}}

% Named special functions.
\newcommand{\Polylogarithm}{\operatorname{Li}}
\newcommand{\LambertW}{W}
\newcommand{\GammaFunction}{\Gamma}
\newcommand{\RiemannZeta}{\zeta}

% Asymptotics.  BigO and LittleO are operator tokens for existing formula
% grammar; prefer the At forms so that the limiting regime is printed.
\newcommand{\BigO}{\mathcal{O}}
\newcommand{\LittleO}{o}
\newcommand{\BigOAt}[2]{\mathcal{O}_{#2}\!\left(#1\right)}
\newcommand{\LittleOAt}[2]{o_{#2}\!\left(#1\right)}

\usepackage{bm}
\usepackage{booktabs,longtable,array,tabularx}
\usepackage{graphicx}
\usepackage{pgfplots}
\usepackage{xcolor}
\usepackage{enumitem}
\usepackage{fancyhdr}
\usepackage{titlesec}
\usepackage{listings}
\usepackage{xurl}
\usepackage{hyperref}
\usepackage[nameinlink,noabbrev]{cleveref}

\definecolor{linkblue}{RGB}{25,75,135}
\definecolor{shadegray}{RGB}{246,247,249}
\definecolor{rulegray}{RGB}{195,201,208}
\pgfplotsset{compat=1.18}
\hypersetup{
  colorlinks=true,
  linkcolor=linkblue,
  citecolor=linkblue,
  urlcolor=linkblue,
  pdftitle={Inverse q-Analogues: Branch Geometry, Asymptotic Inversion, and Computation},
  pdfsubject={Inverse functions of q-Pochhammer symbols, Gaussian coefficients, q-gamma, q-beta, and related q-special functions},
  pdfkeywords={inverse q-Pochhammer, inverse q-binomial, inverse q-gamma, branch points, discriminant, Lagrange inversion, q-series}
}

\pagestyle{fancy}
\fancyhf{}
\fancyhead[L]{\small Inverse \(q\)-analogues}
\fancyhead[R]{\small\nouppercase{\leftmark}}
\fancyfoot[C]{\thepage}
\renewcommand{\headrulewidth}{0.4pt}

\titleformat{\section}{\Large\bfseries}{\thesection}{0.7em}{}
\titleformat{\subsection}{\large\bfseries}{\thesubsection}{0.7em}{}
\titleformat{\subsubsection}{\normalsize\bfseries}{\thesubsubsection}{0.7em}{}
\setcounter{secnumdepth}{3}
\setcounter{tocdepth}{2}
\setlist[itemize]{leftmargin=2em,itemsep=0.25em,topsep=0.4em}
\setlist[enumerate]{leftmargin=2.3em,itemsep=0.25em,topsep=0.4em}

% Long unbreakable inline mathematics occasionally leaves a line with no legal
% break point.  A small emergency stretch lets TeX loosen such a line rather
% than let it protrude into the margin; it applies only to lines that would
% otherwise be overfull.
\setlength{\emergencystretch}{2em}

\newtheorem{theorem}{Theorem}[section]
\newtheorem{proposition}[theorem]{Proposition}
\newtheorem{lemma}[theorem]{Lemma}
\newtheorem{corollary}[theorem]{Corollary}
\newtheorem{conjecture}[theorem]{Conjecture}
\theoremstyle{definition}
\newtheorem{definition}[theorem]{Definition}
\newtheorem{algorithm}[theorem]{Algorithm}
\newtheorem{example}[theorem]{Example}
\theoremstyle{remark}
\newtheorem{remark}[theorem]{Remark}
\newtheorem{warning}[theorem]{Warning}

\newcommand{\lean}[1]{%
  \ifmmode\text{\nolinkurl{#1}}\else\nolinkurl{#1}\fi}

\lstdefinestyle{pseudo}{
  basicstyle=\ttfamily\small,
  backgroundcolor=\color{shadegray},
  frame=single,
  rulecolor=\color{rulegray},
  columns=fullflexible,
  keepspaces=true,
  showstringspaces=false,
  xleftmargin=0.7em,
  xrightmargin=0.7em,
  aboveskip=0.8em,
  belowskip=0.8em
}

% Report-local theorem environments, notation, and status colours.
\theoremstyle{plain}
\newtheorem{principle}[theorem]{Principle}
\newtheorem{identity}[theorem]{Identity}
\theoremstyle{definition}
\newtheorem{researchproblem}[theorem]{Research problem}
\theoremstyle{remark}
\newtheorem{computationalresult}[theorem]{Computational result}

\definecolor{provedgreen}{RGB}{20,100,60}
\definecolor{conjectureorange}{RGB}{158,82,0}
\definecolor{numericalblue}{RGB}{28,78,135}
\newcommand{\Disc}{\operatorname{Disc}}
\newcommand{\Aq}{\mathcal A_q}
\newcommand{\statusproved}{\textcolor{provedgreen}{\textsc{paper proof}}}
\newcommand{\statusnumerical}{\textcolor{numericalblue}{\textsc{symbolic/numerical}}}
\newcommand{\statusconjectural}{\textcolor{conjectureorange}{\textsc{conjectural}}}
\allowdisplaybreaks[3]
\numberwithin{equation}{section}
\numberwithin{figure}{section}
\numberwithin{table}{section}

% The arrival used a book-style three-level hierarchy.  These aliases map it
% onto the canonical article counters without changing mathematical content.
\let\reporttopsection\section
\let\reportsubsection\subsection
\let\reportsubsubsection\subsubsection
\let\chapter\reporttopsection
\let\section\reportsubsection
\let\subsection\reportsubsubsection
\providecommand{\frontmatter}{}
\providecommand{\mainmatter}{}
\providecommand{\backmatter}{}

\title{\bfseries Inverse \(q\)-Analogues\\[0.4em]
\large Branch Geometry, Asymptotic Inversion, and Computation for\\
\(q\)-Pochhammer Symbols, Gaussian Coefficients, and Related Functions}
\author{A research companion to the ProveIt \(q\)-series monograph}
\date{Arrival 30 August 2026; audited against the live repository corpus}

\begin{document}
\frontmatter
\hypersetup{pageanchor=false}
\maketitle
\clearpage
\hypersetup{pageanchor=true}

\chapter*{Abstract}
\addcontentsline{toc}{section}{Abstract}
This report develops a systematic theory of inverse functions for the principal
objects of elementary and basic hypergeometric analysis.  The central point is
that ``the inverse'' is almost never a single globally defined function.  One
must specify the parameter being recovered, the real or complex domain, a
branch, and the singular set across which continuation fails.  We therefore
study partial inverses with respect to the argument, the base \(q\), the order,
and the discrete or continuously interpolated parameters.

For finite \(q\)-Pochhammer products we give a complete real branch
classification in the argument, compute the exact discriminant at level zero,
and introduce a \emph{secondary discriminant} that detects collisions of
critical values.  The first nontrivial positive-base collision occurs already
for \((a;q)_5\): an exact reciprocal degree-ten polynomial has a unique root
\(q_0\in(0,1)\), numerically
\[
 q_0=0.8401533783083766586629364318891329431\ldots .
\]
At this base two distinct negative critical values coincide.  This creates a
Maxwell-type locus in the inverse covering and shows why naive global ordering
conjectures fail.

For the infinite product \((a;q)_\infty\), we define and analyze its principal
real inverse \(\Aq:(0,\infty)\to(-\infty,1)\), derive endpoint series, a
log-periodically corrected large-target asymptotic, and the asymptotic location
and height of every nonprincipal real branch point.  Inversion with respect to
the order is reduced exactly to \(\Aq\).  Inversion with respect to the base is
proved globally monotone in the principal sign regimes and is expanded both at
\(q=0\) and \(q=1\), the latter in polylogarithmic form.

For Gaussian coefficients, positivity of coefficients yields a global theorem:
for every nontrivial \(\GaussianBinomial{n}{k}{q}\) and every \(y>1\), there is exactly one
positive real base \(q\) with \(\GaussianBinomial{n}{k}{q}=y\).  Bernoulli expansions are then
reverted to all orders near \(q=1\), and reciprocal symmetry controls the
large-\(q\) inverse.  Continuous inverses in the upper and lower parameters are
developed through the \(q\)-gamma representation.  Separate chapters treat
inverse \(q\)-numbers, \(q\)-factorials, \(q\)-gamma, \(q\)-digamma,
\(q\)-beta, Euler \(q\)-exponentials, basic hypergeometric series, theta
functions, and broader positive-coefficient \(q\)-analogs.

The report distinguishes proved statements, reproducible symbolic or numerical
findings, and conjectures.  The accompanying Python program regenerates all
figures and data, including high-precision branch-point calculations and exact
secondary discriminants through order six.

\chapter*{Relation to the source monograph and status of claims}
\addcontentsline{toc}{section}{Relation to the source monograph and status of claims}
The source document is Vladimir Reshetnikov's proof-oriented monograph
\emph{\(q\)-Pochhammer Symbols and \(q\)-Binomial Coefficients} in the ProveIt
repository \cite{ReshetnikovMonograph}.  It supplies the forward theory used
here: finite and infinite products, logarithmic derivatives, complex orders,
Gaussian coefficients, \(q\)-gamma and \(q\)-beta functions, and Bernoulli
asymptotics.  The present report is a companion organized around inversion,
branch geometry, conditioning, and algorithms rather than a repetition of that
forward theory.

A label at the beginning of a statement has the following meaning.
\begin{center}
\begin{tabularx}{\textwidth}{@{}l>{\raggedright\arraybackslash}X@{}}
\toprule
\statusproved & a proof is supplied here, or a standard cited theorem is used explicitly;\\
\statusnumerical & exact symbolic algebra or high-precision numerics are reproducible from the archive;\\
\statusconjectural & evidence and a precise formulation are given, but no proof is claimed.\\
\bottomrule
\end{tabularx}
\end{center}
Here and below, ``paper proof'' means proved in this report or imported from a
cited mathematical source; it does \emph{not} assert Lean formalization.  The
live Lean corpus supplies selected finite Gaussian/Pochhammer algebra, including
\lean{Fabius.finiteQPochhammerIn}, \lean{Fabius.gaussianBinomial}, and
\lean{Fabius.finite_qBinomial_theorem}, together with real and complex
infinite-product convergence such as \lean{Fabius.tendsto_finiteQPochhammerIn}
and \lean{Fabius.tendsto_finiteQPochhammerIn_complex}.  The exact base-reversal
crosswalk \lean{Fabius.finiteQPochhammerIn_base_reversal} requires
\(a\ne0\) and \(q\ne0\).  None of this report's specialized inverse,
critical-atlas, discriminant/Maxwell, monodromy, branch-geometry, or inverse
asymptotic theorems currently has an exact public Lean counterpart.  The source
monograph's 184-row formalization ledger records 35 exact, 24 partial, 122
without counterparts, and 3 interface-only rows; in particular, evaluated
Gaussian positivity and infinite-product convergence do not establish the
coefficientwise, derivative-sign, inverse, or asymptotic claims made here.

Some results, especially the secondary-discriminant calculations and the
asymptotic organization of the nonprincipal inverse branches, appear not to be
standard in this form.  No assertion of bibliographic priority is made without
an exhaustive literature search.  Numerical experiments are not substitutes
for interval or formal proof; the exact polynomial identities, however, can be
checked in a computer algebra system over \(\IntegerNumbers[q]\).

\tableofcontents
\listoffigures

\mainmatter

\part{Foundations of parameter-wise inversion}

\chapter{What it means to invert a \texorpdfstring{\(q\)}{q}-analogue}

\section{Partial inverses, not one universal inverse}
Let
\[
 F(p_1,\ldots,p_m)=y
\]
be a scalar \(q\)-analogue with several parameters.  One scalar equation does
not determine all \(m\) parameters: generically its solution set is an
\((m-1)\)-dimensional level hypersurface.  Thus throughout this report an
\emph{inverse with respect to a parameter} means that all other parameters are
fixed and one chosen parameter is recovered.  A simultaneous inverse requires
an equal number of independent equations or additional normalization
conditions.

\begin{definition}[Partial inverse branch]
Fix auxiliary parameters \(\lambda\) and write \(F(p;\lambda)\).  A partial
inverse branch is a function \(p=G(y;\lambda)\) on a domain \(U\) such that
\[
 F(G(y;\lambda);\lambda)=y,
\]
with a specified initial value, real interval, or analytic-continuation path.
Its branch locus consists of singular values reached from points where
\(F_p=0\), together with values caused by poles or singular auxiliary
parameters.
\end{definition}

The main objects and their inversion variables are summarized in
\Cref{tab:objects-variables}.

\begin{longtable}{@{}>{\raggedright\arraybackslash}p{0.24\textwidth}
                         >{\raggedright\arraybackslash}p{0.29\textwidth}
                         >{\raggedright\arraybackslash}p{0.39\textwidth}@{}}
\caption{Inverse variables and principal tools.}\label{tab:objects-variables}\\
\toprule
Object & Recoverable parameter & Principal mechanism\\
\midrule
\endfirsthead
\toprule
Object & Recoverable parameter & Principal mechanism\\
\midrule
\endhead
\((a;q)_n\), \((a;q)_\infty\) & argument \(a\) & algebraic covering for finite \(n\); entire-function covering for \(n=\infty\)\\
& base \(q\) & monotonicity on real sign domains; polynomial or analytic branch continuation\\
& order \(n\) or \(\alpha\) & discrete recurrence or complex-order quotient\\[0.2em]
\(\GaussianBinomial{n}{k}{q}\) & base \(q\) & nonnegative coefficients, Bernoulli reversion, reciprocal symmetry\\
& \(n,k\) or \(\alpha,\beta\) & ratio recurrences and \(q\)-gamma interpolation\\[0.2em]
\([x]_q\), \([n]_q!\) & \(x,n,q\) & elementary logarithm or positive-polynomial inversion\\[0.2em]
\(\Gamma_q(x)\) & \(x\) & two real branches separated by the unique minimum\\
& \(q\) & known monotonicity sectors and endpoint asymptotics\\[0.2em]
\(B_q(a,b)\) & \(a,b\) & strict coordinate monotonicity\\
& \(q\) & Lambert-series derivative; complete integer-parameter theorem\\[0.2em]
Euler \(q\)-exponentials & argument & exact reduction to the inverse infinite product\\[0.2em]
Basic hypergeometric and theta functions & argument and shape parameters & local Lagrange inversion, critical values, quasi-periodicity\\
\bottomrule
\end{longtable}

\section{Real inverse, complex inverse, and Riemann surface}
On a real interval, a continuous strictly monotone function has one ordinary
inverse.  In the complex plane, even a polynomial of degree two has two sheets.
A useful global object is therefore the analytic curve
\[
 \mathscr R_F=\{(p,y)\in\ComplexNumbers^2:F(p)=y\},
\]
projected to the \(y\)-plane.  Away from critical values this projection is a
covering map.  Analytic continuation along loops permutes the sheets; this is
the monodromy action.  For finite \(q\)-Pochhammer products the covering is
algebraic, whereas for the infinite product it has infinitely many sheets and
branch points.

\section{Normalization and branch labels}
A branch can be labelled in several non-equivalent ways:
\begin{enumerate}[label=(\roman*)]
\item by the real interval containing the recovered parameter;
\item by a base point \((p_0,y_0)\) and analytic continuation from it;
\item by an asymptotic condition, such as \(p(y)\to1\) as \(y\to0^+\);
\item by a logarithm index or winding number;
\item by ordering roots at a chosen regular target.
\end{enumerate}
The report uses interval labels for real branches and continuation labels for
complex branches.  A numerical routine that omits such a label is incomplete:
Newton's method can jump between sheets.

\chapter{Local analytic inversion and singularity theory}

\section{The inverse function theorem and sensitivities}
\begin{theorem}[\statusproved\ Local partial inverse]
Let \(F(p,\lambda)\) be holomorphic near \((p_0,\lambda_0)\), put
\(y_0=F(p_0,\lambda_0)\), and assume \(F_p(p_0,\lambda_0)\ne0\).  Then there is
a unique holomorphic branch \(p=G(y,\lambda)\) near
\((y_0,\lambda_0)\) satisfying \(F(G(y,\lambda),\lambda)=y\) and
\(G(y_0,\lambda_0)=p_0\).  Its first sensitivities are
\[
 G_y=\frac1{F_p},
 \qquad
 G_{\lambda_j}=-\frac{F_{\lambda_j}}{F_p},
\]
where the derivatives on the right are evaluated at \(p=G(y,\lambda)\).
\end{theorem}
\begin{proof}
The holomorphic implicit-function theorem gives the branch.  Differentiating
\(F(G(y,\lambda),\lambda)=y\), first in \(y\) and then in \(\lambda_j\),
gives the displayed formulas.
\end{proof}

\begin{corollary}[\statusproved\ First inverse derivatives]
For one variable,
\begin{align*}
 G'(y)&=\frac1{F'(p)},\\
 G''(y)&=-\frac{F''(p)}{F'(p)^3},\\
 G'''(y)&=\frac{3F''(p)^2-F'(p)F'''(p)}{F'(p)^5}.
\end{align*}
\end{corollary}

A dimensionless relative condition number is
\[
 \kappa_{p\leftarrow y}
 =\AbsoluteValue{\frac{y}{pF_p(p)}}.
\]
At \(p=0\), or when relative error is inappropriate, the absolute condition
number \(1/\AbsoluteValue{F_p}\) is preferable.  Divergence of either quantity near
\(F_p=0\) is not a numerical artifact: it is the analytic signature of a
branch point.

\section{Lagrange--Buerhmann coefficients}
\begin{theorem}[\statusproved\ Lagrange--Buerhmann inversion]
Let \(F\) be holomorphic near \(p_0\), with \(F'(p_0)\ne0\), and set
\(y_0=F(p_0)\).  The local inverse is
\[
 G(y)=p_0+\sum_{m\ge1}\frac{(y-y_0)^m}{m!}
 \left.
 \frac{\Differential^{m-1}}{\Differential p^{m-1}}
 \left(\frac{p-p_0}{F(p)-F(p_0)}\right)^m
 \right|_{p=p_0}.
\]
The radius of convergence is bounded by the nearest singular value of the
chosen sheet.
\end{theorem}
\begin{proof}
Apply the classical Lagrange inversion theorem to
\(w=p-p_0\) and \(z=y-y_0=w\phi(w)\), where
\(\phi(0)=F'(p_0)\ne0\).  The coefficient formula follows after rewriting
\(1/\phi(w)=(p-p_0)/(F(p)-F(p_0))\).  Analytic continuation of the inverse can
fail only when the direct map or its local inverse becomes singular, hence a
critical or singular value bounds the disk.
\end{proof}

If
\[
 F(p_0+w)=y_0+f_1w+f_2w^2+f_3w^3+\cdots
\]
and
\[
 G(y_0+z)=p_0+g_1z+g_2z^2+g_3z^3+\cdots,
\]
then
\[
 g_1=\frac1{f_1},\qquad
 g_2=-\frac{f_2}{f_1^3},\qquad
 g_3=\frac{2f_2^2}{f_1^5}-\frac{f_3}{f_1^4}.
\]
This three-term template recurs throughout the report.

\section{Simple branch points and Puiseux inversion}
\begin{theorem}[\statusproved\ Universal square-root law]
Suppose
\[
 F'(p_c)=0,\qquad F''(p_c)\ne0,
 \qquad y_c=F(p_c).
\]
Then the two local inverse sheets have Puiseux expansions
\[
 p=p_c\pm\sqrt{\frac{2(y-y_c)}{F''(p_c)}}
 -\frac{F'''(p_c)}{3F''(p_c)^2}(y-y_c)
 +O\bigl((y-y_c)^{3/2}\bigr),
\]
with a consistent choice of square root.
\end{theorem}
\begin{proof}
Write \(u=p-p_c\), \(v=y-y_c\), and expand
\[
 v=\frac{F''(p_c)}2u^2+\frac{F'''(p_c)}6u^3+O(u^4).
\]
Insert \(u=c_1v^{1/2}+c_2v+O(v^{3/2})\).  Comparing the coefficients of
\(v\) and \(v^{3/2}\) gives
\(c_1^2=2/F''(p_c)\) and
\(c_2=-F'''(p_c)/(3F''(p_c)^2)\).
\end{proof}

More generally, if the first nonzero derivative is \(F^{(r)}(p_c)\), the local
inverse has \(r\)-th-root ramification.  This provides the correct local model
for multiple critical points and for coalescing branches.

\section{Discriminants, caustics, and Maxwell sets}
For a polynomial \(F(p;\lambda)\), the critical-value polynomial is
\[
 \Delta(y;\lambda)=\Disc_p(F(p;\lambda)-y).
\]
Its roots in \(y\) are the critical values, counted with multiplicity.  Two
different degeneracies should be distinguished.
\begin{definition}[Caustic and Maxwell locus]
The \emph{caustic} in parameter space is the set where a critical point itself
becomes multiple, so that \(F_p=F_{pp}=0\).  The \emph{Maxwell locus} is the
set where two distinct simple critical points have the same critical value.
For a one-parameter polynomial family both are detected, with multiplicity, by
\[
 \mathfrak M(\lambda)=\Disc_y\Delta(y;\lambda).
\]
\end{definition}
The terminology is useful because the inverse geometry is different.  At a
caustic, local ramification order changes.  At a Maxwell point, two ordinary
square-root branch points project to the same target value.

\section{Several unknown parameters}
If \(F:\ComplexNumbers^m\to\ComplexNumbers^m\), a simultaneous local inverse exists when
\(\det DF\ne0\).  Its differential is
\[
 \Differential\bm p=(DF)^{-1}\Differential\bm y.
\]
For a scalar \(F:\ComplexNumbers^m\to\ComplexNumbers\), no such matrix inverse exists; the tangent space
of the level set is the nullspace of \(\nabla F\).  This observation prevents a
common ambiguity in phrases such as ``invert with respect to \(q\) and all
parameters'': the mathematically meaningful program is to construct every
one-parameter partial inverse and then study coupled systems separately.

\chapter{Derivative kernels for \texorpdfstring{\(q\)}{q}-products}

\section{Finite and infinite argument derivatives}
For \(N\in\NaturalNumbers\cup\{\infty\}\), write
\[
 P_N(a;q)=(a;q)_N.
\]
For \(N=\infty\) assume \(0<\AbsoluteValue q<1\).  Define
\[
 S_{r,N}(a;q)=\sum_{j=0}^{N-1}
 \frac{q^{rj}}{(1-aq^j)^r},
\]
with the evident infinite interpretation.

\begin{proposition}[\statusproved\ Logarithmic derivative hierarchy]
Away from the zeros of \(P_N\), for every \(r\ge1\),
\[
 \frac{\partial^r}{\partial a^r}\log P_N(a;q)
 =-(r-1)!\,S_{r,N}(a;q).
\]
In particular,
\[
 P_N'(a;q)=-P_N(a;q)S_{1,N}(a;q)
\]
and
\[
 \frac{P_N''(a;q)}{P_N(a;q)}
 =S_{1,N}(a;q)^2-S_{2,N}(a;q).
\]
\end{proposition}
\begin{proof}
Differentiate \(\sum_j\log(1-aq^j)\) term by term.  In the infinite case the
series and all its derivatives converge locally uniformly away from the zero
set.  The final identities follow from differentiating an exponential.
\end{proof}

On \(a<1\) with \(0<q<1\), all factors are positive and
\[
 (-1)^rP_N^{(r)}(a;q)>0
\]
for every available derivative \(1\le r\le N\), and for every \(r\ge1\) when
\(N=\infty\).  For finite \(N\), this follows by differentiating the product
and summing over \(r\)-element subsets; passage to the infinite product is by
locally uniform convergence.

\section{Base derivatives}
\begin{proposition}[\statusproved\ Derivative with respect to the base]
For finite \(N\),
\[
 \frac{\partial}{\partial q}\log(a;q)_N
 =-a\sum_{j=1}^{N-1}\frac{j q^{j-1}}{1-aq^j}.
\]
For \(N=\infty\) and \(\AbsoluteValue q<1\), the same formula holds with the sum extended
to infinity, locally uniformly away from zeros.
\end{proposition}
\begin{proof}
The \(j=0\) factor is independent of \(q\).  Differentiate every other factor.
The infinite derivative series converges normally on compact subsets on which
the denominators stay separated from zero.
\end{proof}

This kernel immediately settles the principal real sign regimes.  If
\(0<a<1\), the derivative is negative; if \(a<0\), it is positive.  When
\(a>1\), denominators change sign as \(q\) crosses the real zeros
\(q=a^{-1/j}\), and a branch-by-branch analysis is unavoidable.

\section{Complex order derivatives}
For a fixed branch of \(q^\alpha\), define
\[
 (a;q)_\alpha=\frac{(a;q)_\infty}{(aq^\alpha;q)_\infty}.
\]
Put \(b=aq^\alpha\).  Then
\begin{align*}
 \frac{\partial}{\partial\alpha}\log(a;q)_\alpha
 &=(\log q)\sum_{j\ge0}\frac{bq^j}{1-bq^j},\\
 \frac{\partial^2}{\partial\alpha^2}\log(a;q)_\alpha
 &=(\log q)^2\sum_{j\ge0}\frac{bq^j}{(1-bq^j)^2}.
\end{align*}
For real \(0<q<1\) and \(0<b<1\), the first quantity is negative and the
second positive.  Thus the order dependence is strictly decreasing and
log-convex on that real branch.

\section{Truncation bounds for certified evaluation}
\begin{lemma}[\statusproved\ Infinite-product tail bound]\label{thm:product-tail}
Let \(0<\AbsoluteValue q<1\), choose \(N\) so that
\(\rho=\AbsoluteValue a\AbsoluteValue q^N<1\), and write
\(T_N=(aq^N;q)_\infty\).  For the logarithm obtained by continuation from
\(a=0\),
\[
 \AbsoluteValue{\log T_N}
 \le \frac{\rho}{(1-\AbsoluteValue q)(1-\rho)}.
\]
Consequently,
\[
 \AbsoluteValue{T_N-1}
 \le \exp\!\left(\frac{\rho}{(1-\AbsoluteValue q)(1-\rho)}\right)-1.
\]
\end{lemma}
\begin{proof}
For \(\AbsoluteValue z<1\), \(\AbsoluteValue{\log(1-z)}\le\AbsoluteValue z/(1-\AbsoluteValue z)\).  Sum this estimate
for \(z=aq^j\), \(j\ge N\), and bound every denominator below by
\(1-\rho\).  The second estimate follows from \(T_N=\exp(\log T_N)\).
\end{proof}

These bounds give a direct route to interval or ball arithmetic.  They also
explain why evaluation in the logarithmic domain is preferable near \(q=1\),
where direct multiplication underflows long before the mathematical product is
zero.

\part{Inverse \texorpdfstring{\(q\)}{q}-Pochhammer theory}

\chapter{Finite products: inversion in the argument}

\section{The real branch decomposition}
Fix \(n\ge1\) and \(0<q<1\), and put
\[
 P_n(a)=\prod_{j=0}^{n-1}(1-aq^j).
\]
Its zeros
\[
 r_j=q^{-j},\qquad 0\le j\le n-1,
\]
are real, simple, and strictly increasing.

\begin{theorem}[\statusproved\ Complete real critical-point structure]
For \(n\ge2\), the following statements hold.
\begin{enumerate}[label=(\roman*)]
\item On \(( -\infty,1)\), \(P_n\) is positive and strictly decreasing from
\(+\infty\) to \(0\).  Hence every \(y>0\) has exactly one inverse there.
\item For each \(k=0,\ldots,n-2\), the interval
\[
 I_k=(q^{-k},q^{-(k+1)})
\]
contains exactly one critical point \(c_{n,k}(q)\).  It is simple.  The sign of
\(P_n\) on \(I_k\) is \((-1)^{k+1}\), and \(P_n(c_{n,k})\) is the unique
extremum on that interval.
\item On the last interval \((q^{-(n-1)},\infty)\), there is no critical point.
The map is monotone from \(0\) to \(+\infty\) if \(n\) is even and from \(0\)
to \(-\infty\) if \(n\) is odd.
\end{enumerate}
\end{theorem}
\begin{proof}
On \(a<1\), every factor is positive and
\[
 \frac{P_n'(a)}{P_n(a)}
 =-S_{1,n}(a;q)<0.
\]
The endpoint limits follow from the leading term
\((-1)^nq^{n(n-1)/2}a^n\) and the zero at \(a=1\).

On \(I_k\), the sum
\[
 S_{1,n}(a;q)=\sum_{j=0}^{n-1}\frac{q^j}{1-aq^j}
\]
is strictly increasing because
\[
 \frac{\partial S_{1,n}}{\partial a}=S_{2,n}>0.
\]
As \(a\downarrow q^{-k}\), the \(j=k\) summand tends to \(-\infty\), while as
\(a\uparrow q^{-(k+1)}\), the \(j=k+1\) summand tends to \(+\infty\).  Thus
there is exactly one zero of \(S_{1,n}\), hence one critical point.  At that
point \(S_{2,n}>0\), so differentiating \(P_n'=-P_nS_{1,n}\) shows
\(P_n''=-P_nS_{2,n}\ne0\); the critical point is simple.  Exactly \(k+1\)
factors are negative on \(I_k\), giving the sign.

On the last interval, \(S_{1,n}\) rises from \(-\infty\) to \(0^-\), so it
never vanishes.  The sign and endpoint behavior give the final assertion.
\end{proof}

\begin{corollary}[\statusproved\ Real root count at a regular target]
Let \(v_{n,k}=P_n(c_{n,k})\).  Every interior interval \(I_k\) contributes two
real inverse values when the target lies strictly between \(0\) and
\(v_{n,k}\), one double inverse value when \(y=v_{n,k}\), and none otherwise.
The principal interval contributes one value for \(y>0\); the final interval
contributes one value whenever \((-1)^ny>0\).  These contributions are
independent and exhaust all real roots of \(P_n(a)=y\).
\end{corollary}

This is a complete real atlas.  In particular, ``the inverse finite
\(q\)-Pochhammer symbol'' is naturally a collection of branches indexed by the
zero intervals.

\section{Exact low-degree inverses}
For \(n=1\), \((a;q)_1=1-a\), so \(a=1-y\).  For \(n=2\),
\[
 (a;q)_2=1-(1+q)a+qa^2,
\]
and therefore
\begin{equation}
 a_{\pm}(y;q)=
 \frac{1+q\pm\sqrt{(1-q)^2+4qy}}{2q}.
 \label{eq:n2-inverse}
\end{equation}
The discriminant in \(y\) vanishes at
\(- (1-q)^2/(4q)\), the unique critical value.  Formula
\eqref{eq:n2-inverse} is the model for every simple branch point: two sheets
exchange when the target circles the critical value.

Whenever the hypotheses of the generic symmetric-monodromy proposition below
hold, the monodromy is \(S_n\); for \(n\ge5\), that branch problem is not
solvable by radicals.  The uniform generic-Galois assertion remains
conjectural.  This is not a deficiency of the theory: discriminants, monodromy,
convergent local series, and certified continuation provide the appropriate
replacement.

\section{The principal series at \texorpdfstring{\(y=1\)}{y=1}}
Euler's finite \(q\)-binomial theorem gives
\[
 P_n(a)=\sum_{m=0}^n(-1)^m q^{\binom m2}\GaussianBinomial{n}{m}{q} a^m.
\]
Set
\[
 e_1=[n]_q,
 \qquad e_2=q\GaussianBinomial{n}{2}{q},
 \qquad e_3=q^3\GaussianBinomial{n}{3}{q},
 \qquad u=1-y.
\]

\begin{proposition}[\statusproved\ Local inverse at the origin]
The branch with \(a(1)=0\) has expansion
\begin{equation}
 a=\frac{u}{e_1}
 +\frac{e_2}{e_1^3}u^2
 +\left(\frac{2e_2^2}{e_1^5}-\frac{e_3}{e_1^4}\right)u^3
 +O(u^4).
 \label{eq:finite-principal-series}
\end{equation}
All higher coefficients are obtained recursively from the finite
\(q\)-binomial coefficients, or directly by Lagrange--Buerhmann inversion.
\end{proposition}
\begin{proof}
The direct series is
\(u=e_1a-e_2a^2+e_3a^3+O(a^4)\).  Substitute
\(a=A_1u+A_2u^2+A_3u^3+O(u^4)\) and compare coefficients.
\end{proof}

Because the nearest singularity of this branch is a critical value, the exact
radius of convergence is the distance from \(1\) to the closest critical value
that is visible on the analytic continuation of the sheet.  Determining that
visibility is a monodromy problem, not merely a nearest-point calculation.

\section{The discriminant of the product}
\begin{theorem}[\statusproved\ Exact zero-level discriminant]
For \(n\ge1\),
\begin{equation}
 \Disc_a P_n(a;q)
 =q^{n(n-1)(n-2)/3}
 \prod_{d=1}^{n-1}(1-q^d)^{2(n-d)}.
 \label{eq:finite-poch-disc}
\end{equation}
\end{theorem}
\begin{proof}
The leading coefficient is
\(L=(-1)^nq^{N}\), \(N=n(n-1)/2\), and the roots are \(q^{-j}\).  Hence
\[
 \Disc P_n=L^{2n-2}\prod_{0\le i<j<n}(q^{-i}-q^{-j})^2.
\]
Now
\[
 q^{-i}-q^{-j}=-q^{-j}(1-q^{j-i}).
\]
For each difference \(d=j-i\), the factor \((1-q^d)^2\) occurs \(n-d\)
times.  The total exponent of \(q\) is
\[
 n(n-1)^2-2\sum_{j=1}^{n-1}j^2
 =\frac{n(n-1)(n-2)}3.
\]
\end{proof}

For a variable target define
\[
 \Delta_n(y;q)=\Disc_a(P_n(a;q)-y).
\]

\begin{proposition}[\statusproved\ Critical-value factorization]
Let \(c_1,\ldots,c_{n-1}\) be the roots of \(P_n'\), counted with
multiplicity.  Then
\begin{equation}
 \Delta_n(y;q)=(-1)^{n(n-1)/2}n^nL^{n-1}
 \prod_{j=1}^{n-1}\bigl(P_n(c_j;q)-y\bigr),
 \qquad L=(-1)^nq^{n(n-1)/2}.
 \label{eq:critical-value-factorization}
\end{equation}
\end{proposition}
\begin{proof}
Use
\[
 \Disc(F)=(-1)^{n(n-1)/2}L^{-1}\ResidueOperator(F,F')
\]
and evaluate the resultant through the roots of \(F'=P_n'\), whose leading
coefficient is \(nL\).
\end{proof}

Equation \eqref{eq:critical-value-factorization} is the algebraic bridge between
forward critical-point calculations and inverse branch geometry.

\section{Secondary discriminants and the first Maxwell collision}
Define
\[
 \mathfrak M_n(q)=\Disc_y\Delta_n(y;q).
\]
For \(0<q<1\), all critical points are simple by the real branch theorem.
Therefore \(\mathfrak M_n(q)=0\) in this interval if and only if two distinct
critical points have the same critical value.

Exact symbolic factorization gives
\begin{align}
 \mathfrak M_3(q)
 &=16q^6(q-1)^6(q^2+q+1)^3,\label{eq:M3}\\
 \mathfrak M_4(q)
 &=16q^{42}(q-1)^{26}(q+1)^6(q^2+1)^8\notag\\
 &\quad\times(9q^4+31q^3+45q^2+31q+9)^3,\label{eq:M4}
\end{align}
and
\begin{align}
 \mathfrak M_5(q)
 &=256q^{150}(q-1)^{60}(q+1)^{14}(q^2+q+1)^4
 (q^4+q^3+q^2+q+1)^{15}\notag\\
 &\quad\times
 (36q^6+97q^5+181q^4+213q^3+181q^2+97q+36)^3\notag\\
 &\quad\times H_5(q)^2,\label{eq:M5}
\end{align}
where
\begin{equation}
\begin{split}
 H_5(q)={}&8q^{10}-5q^9-7q^8-6q^7-8q^6+35q^5\\
          &-8q^4-6q^3-7q^2-5q+8.
\end{split}
\label{eq:H5}
\end{equation}

\begin{theorem}[\statusproved\ First positive-base Maxwell parameter]
The polynomial \(H_5\) has exactly one zero \(q_0\) in \((0,1)\).  At this
parameter the polynomial map \(a\mapsto(a;q_0)_5\) has two distinct simple
critical points with a common critical value.  No such collision occurs for
\(n=3\) or \(n=4\) and \(0<q<1\).
\end{theorem}
\begin{proof}
The factors in \eqref{eq:M3} and \eqref{eq:M4} are nonzero on \((0,1)\), so the
first assertion for \(n=3,4\) follows.

The polynomial \(H_5\) is reciprocal.  Put \(x=q+q^{-1}\).  Dividing by
\(q^5\) gives
\[
 \frac{H_5(q)}{q^5}=h(x)
 =8x^5-5x^4-47x^3+14x^2+53x+37.
\]
For \(0<q<1\), one has \(x>2\).  Now \(h(2)=-1\), while \(h(x)\to+\infty\).
Moreover,
\[
 h'(x)=40x^4-20x^3-141x^2+28x+53,
\]
\[
 h''(x)=160x^3-60x^2-282x+28,
\qquad
 h'''(x)=480x^2-120x-282.
\]
At \(x=2\), all three are positive, and
\(h^{(4)}(x)=960x-120>0\) for \(x\ge2\).  It follows successively that
\(h'''\), \(h''\), \(h'\), and then \(h\) are strictly increasing on
\([2,\infty)\).  Thus \(h\) has exactly one root there, which corresponds to
one reciprocal pair of positive roots of \(H_5\), exactly one lying in
\((0,1)\).

At this root all the other factors in \eqref{eq:M5} are nonzero.  Since the
critical points are simple for positive \(q<1\), vanishing of the secondary
discriminant means that two distinct critical values coincide.
\end{proof}

\begin{computationalresult}[\statusnumerical\ Identification of the colliding sheets]
High-precision computation gives
\begin{align*}
 q_0&=0.84015337830837665866293643188913294310139258\ldots,\\
 c_{5,0}(q_0)&=1.0714044413496100275076056970032887643\ldots,\\
 c_{5,2}(q_0)&=1.5701685824623346994394512430684932918\ldots,
\end{align*}
and
\[
 P_5(c_{5,0};q_0)=P_5(c_{5,2};q_0)
 =-0.0002954183213260369789427474653292421974\ldots .
\]
The residual in the archived 60-digit calculation is below \(4\times10^{-50}\).
The exact factorization proves the existence and uniqueness of the positive
collision parameter; the sheet identification is reproducible numerical
evidence and can be upgraded to a formal interval proof.
\end{computationalresult}

\begin{figure}[htbp]
\centering
\includegraphics[width=0.92\textwidth]{figures/finite_pochhammer_collision.pdf}
\caption{The first positive-base collision of distinct critical values for a
finite \(q\)-Pochhammer product.  At \(q=q_0\), the extrema in the intervals
\(I_0\) and \(I_2\) lie on the same horizontal level.}
\label{fig:finite-collision}
\end{figure}

For \(n=6\), the exact secondary discriminant has another reciprocal factor
with a root
\[
 q\approx0.77737031523884893092862981385481448,
\]
again corresponding numerically to the collision of the \(I_0\) and \(I_2\)
critical values.  The full factor is recorded in
\path{data/secondary_discriminants.txt}.

\section{Monodromy away from the Maxwell set}
\begin{proposition}[\statusproved\ Generic symmetric monodromy]
Fix \(n\) and a base \(q\) for which all critical points of \(P_n\) are simple
and all critical values are distinct.  Then the monodromy group of
\(P_n(a)=y\) over the regular \(y\)-plane is the full symmetric group \(S_n\).
In particular this holds for every \(q\in(0,1)\) for which
\(\mathfrak M_n(q)\ne0\).
\end{proposition}
\begin{proof}
A small loop around a simple critical value exchanges exactly the two sheets
meeting there, hence gives a transposition.  The polynomial covering is
connected, so the monodromy action is transitive.  A transitive permutation
group generated by transpositions is the full symmetric group: form the graph
whose vertices are sheets and whose edges are the generating transpositions;
transitivity makes the graph connected, and edge transpositions of a connected
graph generate \(S_n\).
\end{proof}

At a Maxwell parameter, a loop around the common critical value can produce a
product of two disjoint transpositions.  Splitting the parameter away from the
Maxwell locus resolves it into two ordinary branch values.  This is the
appropriate local unfolding.

\chapter{The infinite product: inversion in the argument}

\section{The principal inverse}
For \(0<q<1\), the entire function
\[
 P(a)=(a;q)_\infty
\]
has simple zeros \(q^{-j}\), \(j\ge0\).  On the zero-free half-line
\(a<1\), it has a canonical inverse.

\begin{definition}[Principal inverse product]
For \(0<q<1\), define \(\Aq:(0,\infty)\to(-\infty,1)\) by
\[
 (\Aq(y);q)_\infty=y.
\]
\end{definition}

\begin{theorem}[\statusproved\ Global principal branch]
The map \(a\mapsto(a;q)_\infty\) is a real-analytic, strictly decreasing
bijection from \((-\infty,1)\) onto \((0,\infty)\).  Consequently \(\Aq\) is
real analytic, strictly decreasing, and convex.  If
\[
 S_r(a)=\sum_{j\ge0}\frac{q^{rj}}{(1-aq^j)^r},
\]
then
\begin{align}
 \Aq'(y)&=-\frac1{yS_1(\Aq(y))},\label{eq:Aq-prime}\\
 \Aq''(y)&=
 \frac{S_1(\Aq(y))^2-S_2(\Aq(y))}
 {y^2S_1(\Aq(y))^3}>0.\label{eq:Aq-second}
\end{align}
\end{theorem}
\begin{proof}
Every factor is positive for \(a<1\), and the logarithmic derivative is
\(-S_1<0\).  The product tends to zero as \(a\uparrow1\).  As
\(a\to-\infty\), every finite truncation tends to infinity and is bounded above
by the infinite product, so the latter tends to infinity.  The inverse exists.
Formula \eqref{eq:Aq-prime} follows from implicit differentiation.  Since
\[
 S_1^2-S_2=2\sum_{0\le i<j}\frac{q^{i+j}}
 {(1-aq^i)(1-aq^j)}>0,
\]
formula \eqref{eq:Aq-second} proves convexity.
\end{proof}

\section{Three distinguished target regimes}
\subsection{Near \(y=1\)}
The infinite \(q\)-binomial theorem gives
\[
 (a;q)_\infty
 =1-\frac{a}{1-q}
 +\frac{q a^2}{(1-q)(1-q^2)}
 -\frac{q^3a^3}{(1-q)(1-q^2)(1-q^3)}+\cdots .
\]

\begin{proposition}[\statusproved\ Central principal series]
With \(u=1-y\),
\begin{equation}
\begin{split}
 \Aq(y)={}&(1-q)u
 +\frac{q(1-q)}{1+q}u^2\\
 &+\frac{q^2(1-q)(q^2+q+2)}{(1+q)^2(q^2+q+1)}u^3
 +O(u^4).
\end{split}
\label{eq:Aq-series-one}
\end{equation}
\end{proposition}
\begin{proof}
Apply the three-term series-reversion formula to the first three nonconstant
coefficients of the direct product.
\end{proof}

\subsection{Near the first zero, \(y=0^+\)}
Let
\[
 R_q=(q;q)_\infty,
 \qquad
 L_q=\sum_{j=1}^\infty\frac{q^j}{1-q^j}.
\]
Writing \(a=1-\delta\) gives
\[
 (a;q)_\infty
 =\delta(aq;q)_\infty
 =R_q\delta\bigl(1+L_q\delta+O(\delta^2)\bigr).
\]
Hence
\begin{equation}
 \Aq(y)=1-\frac{y}{R_q}
 +L_q\left(\frac{y}{R_q}\right)^2+O(y^3).
 \label{eq:Aq-near-zero}
\end{equation}
This expansion is especially useful when solving Euler \(q\)-exponential
equations at large positive values.

\subsection{Large targets and log-periodic correction}
The behavior as \(y\to\infty\) is not a simple power law.  Put
\(q=\EulerE^{-\lambda}\), \(\lambda>0\), and write \(x=-a>0\),
\(u=\log x\).  Choose the unique decomposition
\[
 x=q^{-m}r,
 \qquad m\in\NaturalNumbers,
 \qquad 1\le r<q^{-1},
\]
and set \(\rho=\log r\in[0,\lambda)\).

\begin{theorem}[\statusproved\ Theta-modulated large-target law]
Define the bounded \(\lambda\)-periodic phase function
\[
 \Omega_q(\rho)=
 \log\bigl((-r;q)_\infty(-q/r;q)_\infty\bigr)
 -\frac{\rho}{2}-\frac{\rho^2}{2\lambda},
 \qquad r=\EulerE^\rho.
\]
Then, as \(x\to\infty\),
\begin{equation}
 \log(-x;q)_\infty
 =\frac{u^2}{2\lambda}+\frac{u}{2}
 +\Omega_q(u\bmod\lambda)+O(\EulerE^{-u}).
 \label{eq:logperiodic-product}
\end{equation}
Consequently, as \(y\to\infty\),
\begin{equation}
 \log\bigl(-\Aq(y)\bigr)
 =\sqrt{2\lambda\log y}-\frac{\lambda}{2}
 +O\bigl((\log y)^{-1/2}\bigr),
 \label{eq:Aq-large-leading}
\end{equation}
with a computable log-periodic correction at the displayed error scale.
\end{theorem}
\begin{proof}
The exact finite shift is
\begin{align*}
 (-q^{-m}r;q)_\infty
 &=\prod_{j=0}^{m-1}(1+rq^{j-m})(-r;q)_\infty\\
 &=r^m q^{-m(m+1)/2}(-q/r;q)_m(-r;q)_\infty.
\end{align*}
The finite product \((-q/r;q)_m\) differs from its infinite limit by a relative
factor \(1+O(q^m)\), uniformly for \(1\le r\le q^{-1}\).  Taking logarithms and
using \(u=m\lambda+\rho\) yields
\[
 \frac{\lambda m(m+1)}2+m\rho
 =\frac{u^2}{2\lambda}+\frac u2
 -\frac{\rho}{2}-\frac{\rho^2}{2\lambda}.
\]
This proves \eqref{eq:logperiodic-product}.  Complete the square in the first
two terms and invert.  Since \(\Omega_q\) is bounded, it perturbs \(u\) by
\(O((\log y)^{-1/2})\).
\end{proof}

The periodic phase is a genuine lattice effect.  Replacing the product sum by a
single Euler--Maclaurin integral captures the quadratic envelope but misses this
fine structure.

\section{All nonprincipal real branches}
For every \(k\ge0\), let
\[
 I_k=(q^{-k},q^{-(k+1)}).
\]
The same proof as in the finite case gives one simple critical point
\(c_k(q)\in I_k\).  The sign on \(I_k\) is \((-1)^{k+1}\), and both endpoints
map to zero.  Thus a target strictly between zero and the signed extremum has
two inverse values in that interval.

\begin{theorem}[\statusproved\ Asymptotics of remote branch points]
Let \(m=k+1\).  As \(k\to\infty\),
\begin{equation}
 c_k(q)=q^{-m}
 \left(1-\frac1{m+1}+O(m^{-2})\right),
 \label{eq:ck-asymptotic}
\end{equation}
and
\begin{equation}
\begin{split}
 \AbsoluteValue{(c_k(q);q)_\infty}
 ={}&(q;q)_\infty^2q^{-m(m+1)/2}
 \frac{m^m}{(m+1)^{m+1}}\bigl(1+O(m^{-1})\bigr)\\
 ={}&\frac{(q;q)_\infty^2}{\EulerE(m+1)}
 q^{-m(m+1)/2}\bigl(1+O(m^{-1})\bigr).
\end{split}
\label{eq:vk-asymptotic}
\end{equation}
The constants implied by the \(O\)-terms may be chosen locally uniformly for
\(q\) in compact subsets of \((0,1)\).
\end{theorem}
\begin{proof}
Write
\[
 a=q^{-m}(1-\varepsilon),
 \qquad 0<\varepsilon<1-q.
\]
A direct split at the zero index gives the exact identity
\begin{equation}
\begin{split}
 (a;q)_\infty={}&(-1)^m q^{-m(m+1)/2}
 \varepsilon(1-\varepsilon)^mG_m(\varepsilon),\\
 G_m(\varepsilon)={}&
 \left(\frac{q}{1-\varepsilon};q\right)_m
 ((1-\varepsilon)q;q)_\infty.
\end{split}
\label{eq:scaled-critical-product}
\end{equation}
For \(\varepsilon=O(m^{-1})\), the logarithmic derivative of \(G_m\) is
uniformly bounded and
\[
 G_m(\varepsilon)=(q;q)_\infty^2(1+O(m^{-1})).
\]
At the extremum,
\[
 0=\frac1\varepsilon-\frac{m}{1-\varepsilon}
 +\frac{\partial}{\partial\varepsilon}\log G_m(\varepsilon).
\]
The bounded final term perturbs the solution of the first two terms by
\(O(m^{-2})\), so
\(\varepsilon=1/(m+1)+O(m^{-2})\).  Substitution in
\eqref{eq:scaled-critical-product} yields
\eqref{eq:vk-asymptotic}; Stirling's elementary limit
\((1-1/(m+1))^m=\EulerE^{-1}(1+O(m^{-1}))\) gives the second line.
\end{proof}

\begin{figure}[htbp]
\centering
\includegraphics[width=0.88\textwidth]{figures/infinite_critical_asymptotics.pdf}
\caption{Ratio of exact nonprincipal critical-value magnitudes to the leading
asymptotic in \eqref{eq:vk-asymptotic}.  Convergence is slower when \(q\) is
close to one, as expected from the product tail.}
\label{fig:infinite-critical-asymptotics}
\end{figure}

At each \(v_k=(c_k;q)_\infty\), the two local branches obey the universal
square-root law
\[
 a=c_k\pm\sqrt{\frac{2(y-v_k)}{P''(c_k)}}+O(y-v_k).
\]
Equations \eqref{eq:ck-asymptotic}--\eqref{eq:vk-asymptotic} therefore give an
asymptotic atlas of the entire real inverse surface.

\section{A Bernstein-function conjecture}
Let
\[
 X_q(y)=1-\Aq(y),\qquad y>0.
\]
Then \(X_q>0\), \(X_q'>0\), and \(X_q''<0\).  The archived series-reversion
experiment finds
\[
 (-1)^{r-1}X_q^{(r)}(y)>0,
 \qquad 1\le r\le8,
\]
at fifteen representative points with
\(q\in\{0.25,0.5,0.75\}\) and arguments ranging from \(-3\) to \(0.8\).

\begin{conjecture}[\statusconjectural\ Complete Bernstein principal inverse]
For every \(0<q<1\), the function \(X_q=1-\Aq\) is a complete Bernstein
function on \((0,\infty)\).  Equivalently, \(X_q'\) is completely monotone and
\(X_q\) admits a Pick--Nevanlinna continuation with the standard complete
Bernstein representation.
\end{conjecture}

A proof would simultaneously control all inverse derivatives, provide sharp
continued-fraction approximants, and connect the zero sequence \(q^{-j}\) to
Stieltjes moment theory.  The conjecture is substantially stronger than the
proved convexity of \(\Aq\).

\begin{figure}[htbp]
\centering
\includegraphics[width=0.88\textwidth]{figures/principal_inverse.pdf}
\caption{The principal inverse \(\Aq(y)\).  It decreases from \(1\) to
\(-\infty\), crossing zero at \(y=1\).}
\label{fig:principal-inverse}
\end{figure}

\chapter{Inversion with respect to the base \texorpdfstring{\(q\)}{q}}

\section{Finite products on the principal sign domains}
Fix \(a\in\RealNumbers\) and \(n\ge2\).  Regard \((a;q)_n\) as a polynomial in \(q\) of
degree \(n(n-1)/2\).

\begin{theorem}[\statusproved\ Global finite-base inverses]
\begin{enumerate}[label=(\roman*)]
\item If \(0<a<1\), the map \(q\mapsto(a;q)_n\) is strictly decreasing on
\((0,1)\), from \(1-a\) to \((1-a)^n\).  Thus there is exactly one real inverse
\(q\in(0,1)\) for
\[
 (1-a)^n<y<1-a.
\]
\item If \(a<0\), the same map is strictly increasing on \((0,1)\), from
\(1-a\) to \((1-a)^n\).  Thus there is exactly one inverse for
\[
 1-a<y<(1-a)^n.
\]
\item If \(a>1\), the real zeros \(q=a^{-1/j}\), \(1\le j\le n-1\), divide
\((0,1)\) into sign intervals.  On each interval inverse branches are governed
by the zeros of the explicit base derivative kernel; no global monotonicity
claim is valid without specifying the interval.
\end{enumerate}
\end{theorem}
\begin{proof}
For \(0<a<1\), every factor with \(j\ge1\) decreases strictly with \(q\), while
the \(j=0\) factor is fixed and positive.  For \(a<0\), every factor increases.
The endpoint limits are immediate.  For \(a>1\), factors vanish at the stated
bases and the logarithmic derivative has poles there.
\end{proof}

\section{The small-base inverse series}
Normalize
\[
 s=\frac{1-y/(1-a)}{a},
 \qquad a\ne0.
\]
For a target near \(1-a\), one obtains
\[
 s=q+q^2-aq^3+O(q^4),\qquad n=3,
\]
and
\[
 s=q+q^2+(1-a)q^3+O(q^4),\qquad n\ge4.
\]
Therefore
\begin{align}
 q&=s-s^2+(2+a)s^3+O(s^4),&&n=3,\label{eq:qinv-small-n3}\\
 q&=s-s^2+(1+a)s^3+O(s^4),&&n\ge4.\label{eq:qinv-small-n4}
\end{align}
For \(n=2\), the exact answer is \(q=s\).  These formulas also apply to the
infinite product because factors beginning with \(q^4\) do not affect the
cubic truncation.

\section{The infinite product as a function of the base}
\begin{theorem}[\statusproved\ Global infinite-base inverse]
Fix \(a\in\RealNumbers\).
\begin{enumerate}[label=(\roman*)]
\item If \(0<a<1\), then \(q\mapsto(a;q)_\infty\) is strictly decreasing on
\((0,1)\), from \(1-a\) to \(0\).  Every \(y\in(0,1-a)\) has exactly one base
inverse, denoted \(Q_a(y)\).
\item If \(a<0\), the map is strictly increasing on \((0,1)\), from \(1-a\) to
\(+\infty\).  Every \(y>1-a\) has exactly one base inverse.
\item If \(a=0\), the product is identically one.
\end{enumerate}
\end{theorem}
\begin{proof}
The derivative kernel has the asserted sign.  For \(0<a<1\), as \(q\to1^-\),
infinitely many factors are bounded above by a fixed number less than one, so
the product tends to zero.  For \(a<0\), arbitrarily many factors exceed
\(1+\AbsoluteValue a/2\) when \(q\) is sufficiently close to one, so the product tends to
infinity.  The \(q\to0\) limit is \(1-a\).
\end{proof}

\section{Polylogarithmic inversion near \(q=1\)}
Let \(0<a<1\), put \(q=\EulerE^{-t}\), and define
\[
 A=\Polylogarithm_2(a),
 \qquad B=\frac12\log(1-a),
 \qquad C=\frac{a}{12(1-a)}.
\]
The classical expansion of the shifted factorial is
\begin{equation}
 \log(a;\EulerE^{-t})_\infty
 =-\frac{\Polylogarithm_2(a)}t+\frac12\log(1-a)
 -\sum_{r\ge1}\frac{B_{2r}}{(2r)!}
 \Polylogarithm_{2-2r}(a)t^{2r-1}.
 \label{eq:mcintosh-expansion}
\end{equation}
It is asymptotic in sectors and, after truncation, admits standard remainder
bounds; see McIntosh \cite{McIntosh1999} and later Lambert-series treatments
\cite{BanerjeeWilkerson}.

\begin{theorem}[\statusproved\ Base inverse as \(y\to0^+\)]
Let
\[
 X=B-\log y=\log\frac{\sqrt{1-a}}{y}.
\]
Then the unique base inverse satisfies
\begin{equation}
 t=-\log Q_a(y)
 =\frac{A}{X}+\frac{CA^2}{X^3}+O(X^{-5}),
 \qquad y\to0^+,
 \label{eq:qbase-inverse-near1}
\end{equation}
so \(Q_a(y)=\exp(-t)\).  Higher terms follow by formal reversion of
\eqref{eq:mcintosh-expansion}.
\end{theorem}
\begin{proof}
From \eqref{eq:mcintosh-expansion},
\[
 X=\frac At+Ct+O(t^3).
\]
Insert \(t=A/X+d/X^3+O(X^{-5})\).  Expanding \(A/t+Ct\) gives
\[
 X+\frac{-d/A+CA}{X}+O(X^{-3}),
\]
so \(d=CA^2\).
\end{proof}

The dominant inverse scale is rational in \(1/\log(1/y)\), not algebraic in
\(y\).  Joint limits \(a\to1\) and \(q\to1\) can change this scale and create
Lambert-\(W\) transition formulas; they are listed among the frontier problems
in \Cref{chap:research}.

\section{Complex base branches}
For finite \(n\), solving for \(q\) is algebraic of degree
\(n(n-1)/2\), with branch points given by
\[
 \Disc_q((a;q)_n-y)=0.
\]
The reciprocal identity
\begin{equation}
 (a;q)_n=(-a)^nq^{n(n-1)/2}(a^{-1};q^{-1})_n
 \label{eq:poch-reciprocal}
\end{equation}
transports large-base behavior to small reciprocal base, provided zeros and
logarithm branches are tracked.

For the infinite product, \(q\mapsto(a;q)_\infty\) is analytic in \(\AbsoluteValue q<1\).
Its inverse branches can be continued until either
\(\partial_q(a;q)_\infty=0\) or the path meets a singular boundary.  Roots of
unity organize the boundary asymptotics, but a blanket assertion of a natural
boundary depends on the chosen parameter and normalization and is not made
here.

\chapter{Inversion with respect to order}

\section{Exact reduction to the principal argument inverse}
Fix \(0<q<1\), and initially let \(0<a<1\).  Put
\[
 P(a)=(a;q)_\infty,
 \qquad
 (a;q)_\alpha=\frac{P(a)}{P(aq^\alpha)}.
\]
Let
\[
 \alpha_0=\log_q(1/a)<0.
\]

\begin{theorem}[\statusproved\ Real order inverse]
On \((\alpha_0,\infty)\), the function
\(\alpha\mapsto(a;q)_\alpha\) is strictly decreasing and log-convex, from
\(+\infty\) to \(P(a)\).  For every \(y>P(a)\), its unique real inverse is
\begin{equation}
 \alpha=
 \frac{\log\!\left(\Aq(P(a)/y)/a\right)}{\log q}.
 \label{eq:order-inverse-exact}
\end{equation}
On the nonnegative-order subbranch \(\alpha\ge0\), the target range is
\(P(a)<y\le1\); the lower endpoint is approached only as
\(\alpha\to\infty\).
\end{theorem}
\begin{proof}
The derivative formulas in the previous chapter show strict decrease and
log-convexity while \(0<aq^\alpha<1\).  At \(\alpha\downarrow\alpha_0\), the
denominator has its first zero and the quotient diverges.  At infinity the
denominator tends to one.  Solving
\[
 P(aq^\alpha)=\frac{P(a)}y
\]
gives \(aq^\alpha=\Aq(P(a)/y)\), hence
\eqref{eq:order-inverse-exact}.
\end{proof}

For \(a<0\), the same algebraic reduction remains valid on the real branch,
with monotonicity reversed: \((a;q)_\alpha\) increases from \(1\) at
\(\alpha=0\) toward the unattained limit \((a;q)_\infty>1\) as
\(\alpha\to\infty\).

\section{Complex order lattice}
Because \(q^\alpha=\exp(\alpha\log q)\), the complex order is periodic under
\[
 \alpha\longmapsto\alpha+\frac{2\pi\ImaginaryUnit m}{\log q},
 \qquad m\in\IntegerNumbers.
\]
If \(b\) is any chosen inverse value satisfying
\((b;q)_\infty=P(a)/y\), then the order sheets are
\begin{equation}
 \alpha_m=
 \frac{\PrincipalLogarithm(b/a)+2\pi\ImaginaryUnit m}{\log q},
 \qquad m\in\IntegerNumbers.
 \label{eq:complex-order-lattice}
\end{equation}
Thus order inversion factors into two independent branch structures: the
infinite-product inverse selecting \(b\), and the logarithmic lattice selecting
\(m\).

\section{Discrete order recovery}
For \(0<a<1\), the sequence \((a;q)_n\) decreases strictly from \(1\) to
\((a;q)_\infty\).  A target has an exact integer inverse only if it belongs to
this sequence.  For arbitrary targets in the interval, the continuous order
inverse \eqref{eq:order-inverse-exact} supplies a canonical interpolation; its
floor and ceiling bracket the target through
\[
 (a;q)_{\lceil\alpha\rceil}
 \le y\le
 (a;q)_{\lfloor\alpha\rfloor}.
\]
This is preferable to an ad hoc nearest-integer search because it preserves
sensitivity information and extends analytically.

\section{Multiple shifted factors}
For
\[
 F(\alpha)=\prod_{r=1}^s(a_r;q)_\alpha^{\mu_r},
\]
where the exponents \(\mu_r\) are fixed, the logarithmic order derivative is
\[
 \frac{F'(\alpha)}{F(\alpha)}
 =\log q\sum_{r=1}^s\mu_r
 \sum_{j\ge0}\frac{a_rq^{\alpha+j}}{1-a_rq^{\alpha+j}}.
\]
Any sign theorem for the combined Lambert kernel yields a global inverse; in
the absence of such a sign, the zeros of this kernel divide the order line into
monotone branches.  This single formula covers many ratios occurring in basic
hypergeometric coefficients.

\part{Gaussian coefficients and combinatorial \texorpdfstring{\(q\)}{q}-analogs}

\chapter{Inverting a Gaussian coefficient with respect to \texorpdfstring{\(q\)}{q}}

\section{A global positive-base theorem}
For integers \(0\le k\le n\), write
\[
 G_{n,k}(q)=\GaussianBinomial{n}{k}{q}
 =\sum_{j=0}^{D}c_jq^j,
 \qquad D=k(n-k).
\]
The coefficients \(c_j\) are nonnegative integers, \(c_0=c_D=1\), and
\(c_j=c_{D-j}\).

\begin{theorem}[\statusproved\ Unique positive-base inverse]
Let \(0<k<n\).  The map \(q\mapsto G_{n,k}(q)\) is strictly increasing on
\((0,\infty)\), from \(1\) to \(+\infty\).  Hence for every target \(y>1\)
there is exactly one positive real solution, denoted
\[
 q=\mathcal Q_{n,k}(y).
\]
Moreover,
\[
 \begin{cases}
 0<\mathcal Q_{n,k}(y)<1,&1<y<\binom nk,\\
 \mathcal Q_{n,k}(y)=1,&y=\binom nk,\\
 \mathcal Q_{n,k}(y)>1,&y>\binom nk.
 \end{cases}
\]
\end{theorem}
\begin{proof}
Since \(D>0\), at least one coefficient with positive exponent is nonzero, and
\[
 G_{n,k}'(q)=\sum_{j=1}^D jc_jq^{j-1}>0
\]
for \(q>0\).  The endpoint values are \(G(0)=1\),
\(G(1)=\binom nk\), and \(G(q)\sim q^D\) as \(q\to\infty\).
\end{proof}

This theorem is stronger than a local implicit-function result: it gives a
single globally defined positive branch for every nontrivial Gaussian
coefficient and every admissible target.

\section{Expansion at \texorpdfstring{\(q=0\)}{q=0}}
Let
\[
 G_{n,k}(q)=1+q+c_2q^2+c_3q^3+\cdots,
\]
where \(c_j\) counts partitions of \(j\) fitting in a
\(k\times(n-k)\) rectangle.  For \(u=y-1\), series reversion gives
\begin{equation}
 \mathcal Q_{n,k}(y)
 =u-c_2u^2+(2c_2^2-c_3)u^3+O(u^4).
 \label{eq:qbinom-inverse-zero}
\end{equation}
Here \(c_2\) is the number of the partitions \((2)\) and \((1,1)\) that fit
the rectangle.  More generally, every inverse coefficient is an integer
polynomial in the rectangle-partition numbers.  Determining its sign pattern is
an interesting combinatorial reversion problem.

\section{Complete Bernoulli reversion near \(q=1\)}
Put \(q=\EulerE^{-t}\), \(D=k(n-k)\), and
\[
 B=\binom nk,
 \qquad
 s=-\log\frac{y}{B}.
\]
For the \(q<1\) branch, \(s>0\).  The forward expansion is
\begin{equation}
 \log\frac{G_{n,k}(\EulerE^{-t})}{B}
 =-\frac D2t+
 \sum_{r\ge1}\frac{B_{2r}t^{2r}}{2r(2r)!}
 \sum_{j=1}^k\bigl((n-k+j)^{2r}-j^{2r}\bigr).
 \label{eq:gaussian-bernoulli-repeated}
\end{equation}
This is the Bernoulli expansion established in the source monograph and is
convergent for fixed \(n,k\) when \(\AbsoluteValue t<2\pi/n\).

\begin{theorem}[\statusproved\ First inverse terms at \(q=1\)]
As \(y\uparrow B\) from below,
\begin{equation}
 t=-\log\mathcal Q_{n,k}(y)
 =\frac{2s}{D}
 +\frac{n+1}{3D^2}s^2
 +\frac{(n+1)^2}{9D^3}s^3
 +O(s^4).
 \label{eq:qbinom-inverse-one}
\end{equation}
For the branch \(q>1\), put \(r=\log(y/B)>0\); then
\begin{equation}
 \log\mathcal Q_{n,k}(y)
 =\frac{2r}{D}
 -\frac{n+1}{3D^2}r^2
 +\frac{(n+1)^2}{9D^3}r^3
 +O(r^4).
 \label{eq:qbinom-inverse-one-plus}
\end{equation}
\end{theorem}
\begin{proof}
The coefficient of \(t^2\) in
\eqref{eq:gaussian-bernoulli-repeated} is
\[
 \frac1{24}\sum_{j=1}^k((n-k+j)^2-j^2)
 =\frac{D(n+1)}{24}.
\]
Thus
\[
 s=\frac D2t-\frac{D(n+1)}{24}t^2+O(t^4).
\]
There is no cubic term because the coefficient distribution is symmetric about
\(D/2\).  Reversion gives \eqref{eq:qbinom-inverse-one}.  Replacing \(t\) by
\(-t\), or using palindromicity directly, gives
\eqref{eq:qbinom-inverse-one-plus}.
\end{proof}

The fourth inverse coefficient can also be written explicitly.  Define
\[
 K_4=\frac1{2880}
 \sum_{j=1}^k\bigl((n-k+j)^4-j^4\bigr).
\]
If \(A=D/2\) and \(C=D(n+1)/24\), then the coefficient of \(s^4\) in
\(t(s)\) is
\[
 \frac{5C^3}{A^7}-\frac{K_4}{A^5}.
\]
Higher coefficients follow from ordinary Bell-polynomial series reversion.
Thus the forward Bernoulli expansion automatically generates a complete
inverse expansion.

\begin{figure}[htbp]
\centering
\includegraphics[width=0.88\textwidth]{figures/qbinomial_inverse_accuracy.pdf}
\caption{Relative errors of the quadratic and cubic truncations of
\eqref{eq:qbinom-inverse-one} for \(\GaussianBinomial{12}{5}{q}\).  The observed slopes
agree with the predicted next powers.}
\label{fig:qbinom-inverse-accuracy}
\end{figure}

\section{The large-target inverse}
Palindromicity gives
\begin{equation}
 G_{n,k}(q)=q^DG_{n,k}(q^{-1}).
 \label{eq:qbinom-reciprocity}
\end{equation}
Since the coefficient of \(q\) is one,
\[
 G_{n,k}(q)=q^D\bigl(1+q^{-1}+O(q^{-2})\bigr).
\]

\begin{proposition}[\statusproved\ Large-target expansion]
As \(y\to\infty\),
\begin{equation}
 \mathcal Q_{n,k}(y)
 =y^{1/D}-\frac1D+O(y^{-1/D}).
 \label{eq:qbinom-large-y}
\end{equation}
A full expansion in powers of \(y^{-1/D}\) is obtained by reverting
\eqref{eq:qbinom-reciprocity} at \(q^{-1}=0\).
\end{proposition}
\begin{proof}
Put \(Y=y^{1/D}\) and \(q=Y+c+O(Y^{-1})\).  Then
\(q^D=Y^D(1+Dc/Y+O(Y^{-2}))\) and
\(1+q^{-1}=1+Y^{-1}+O(Y^{-2})\).  Equality with \(y=Y^D\) forces
\(Dc+1=0\).
\end{proof}

\section{Complex branches and cyclotomic geometry}
The positive real branch is only one sheet of the degree-\(D\) algebraic
inverse.  The cyclotomic factorization
\begin{equation}
 \GaussianBinomial{n}{k}{q}
 =\prod_{d\ge2}\Phi_d(q)^{
 \lfloor n/d\rfloor-\lfloor k/d\rfloor-\lfloor(n-k)/d\rfloor}
 \label{eq:qbinom-cyclotomic}
\end{equation}
shows how roots of unity enter the zero set.  Critical values are the roots of
\[
 \Disc_q(G_{n,k}(q)-y).
\]
Analytic continuation of \(\mathcal Q_{n,k}\) around those values permutes the
\(D\) complex sheets.  Reciprocal symmetry relates the covering near zero to
the covering near infinity but does not identify equal-target sheets directly,
because the target is rescaled by \(q^D\).

\chapter{Inverting the upper and lower Gaussian parameters}

\section{Discrete parameters \texorpdfstring{\(n\) and \(k\)}{n and k}}
Fix \(q>0\).  The standard ratio identities give
\begin{align}
 \frac{\GaussianBinomial{n+1}{k}{q}}{\GaussianBinomial{n}{k}{q}}
 &=\frac{[n+1]_q}{[n+1-k]_q}>1,
 &&0<k\le n,\label{eq:n-ratio}\\
 \frac{\GaussianBinomial{n}{k+1}{q}}{\GaussianBinomial{n}{k}{q}}
 &=\frac{[n-k]_q}{[k+1]_q}.
 \label{eq:k-ratio}
\end{align}
Because \([x]_q\) is increasing in a positive integer \(x\), the sequence in
\(n\) is strictly increasing for fixed \(k\), while the sequence in \(k\) is
symmetric and strictly increases up to the middle, then decreases.

\begin{theorem}[\statusproved\ Discrete inverse structure]
For fixed \(q>0\) and \(k\ge1\), \(n\mapsto\GaussianBinomial{n}{k}{q}\) is strictly increasing
for \(n\ge k\).  For \(0<q<1\), it converges to \(1/(q;q)_k\); for \(q\ge1\)
it is unbounded (at \(q=1\), this is the usual classical binomial growth).

For fixed \(n\), the map \(k\mapsto\GaussianBinomial{n}{k}{q}\) is symmetric under
\(k\leftrightarrow n-k\), strictly increasing for \(k<(n-1)/2\), and strictly
decreasing after the center.  A target therefore has at most two integer
\(k\)-preimages, paired by symmetry.
\end{theorem}
\begin{proof}
The monotonicity follows from \eqref{eq:n-ratio}--\eqref{eq:k-ratio}.  For
\(0<q<1\), use the product formula and let \(n\to\infty\): each numerator factor
tends to one, leaving \(1/(q;q)_k\).
\end{proof}

This discrete structure is exactly mirrored by the continuous interpolation
below.

\section{A two-parameter analytic interpolation}
For \(\alpha\ge\beta\ge0\), define
\begin{equation}
 \mathcal G_q(\alpha,\beta)
 =\frac{\Gamma_q(\alpha+1)}
 {\Gamma_q(\beta+1)\Gamma_q(\alpha-\beta+1)}.
 \label{eq:generalized-gaussian}
\end{equation}
It agrees with \(\GaussianBinomial{n}{k}{q}\) at nonnegative integers.

\subsection{Recovering the upper parameter}
\begin{theorem}[\statusproved\ Upper-parameter inverse]
Fix \(0<q<1\) and \(\beta>0\).  The function
\(\alpha\mapsto\mathcal G_q(\alpha,\beta)\) is strictly increasing and strictly
log-concave on \([\beta,\infty)\), from \(1\) to
\begin{equation}
 L_{q,\beta}=\frac1{(q;q)_\beta}
 =\frac1{(1-q)^\beta\Gamma_q(\beta+1)}.
 \label{eq:upper-limit}
\end{equation}
Hence every \(y\in(1,L_{q,\beta})\) has exactly one upper-parameter inverse.
\end{theorem}
\begin{proof}
Let \(\psi_q=(\log\Gamma_q)'\).  Then
\[
 \partial_\alpha\log\mathcal G_q
 =\psi_q(\alpha+1)-\psi_q(\alpha-\beta+1)>0
\]
because \(\psi_q\) is strictly increasing.  Its derivative is
\[
 \psi_q'(\alpha+1)-\psi_q'(\alpha-\beta+1)<0
\]
because \(\psi_q'\) is strictly decreasing.  The endpoint at
\(\alpha=\beta\) is one.  The product representation
\[
 \mathcal G_q(\alpha,\beta)
 =\frac{(q^{\beta+1};q)_\infty}{(q;q)_\infty}
 \frac{(q^{\alpha-\beta+1};q)_\infty}
 {(q^{\alpha+1};q)_\infty}
\]
gives the limit.
\end{proof}

Put \(z=q^{\alpha-\beta+1}\).  The exact factorization
\[
 \mathcal G_q(\alpha,\beta)=L_{q,\beta}(z;q)_\beta
\]
reduces the upper-parameter inverse to order or argument inversion.  As
\(y\uparrow L_{q,\beta}\),
\[
 L_{q,\beta}-y
 =L_{q,\beta}[\beta]_qz+O(z^2),
\]
so
\begin{equation}
 \alpha=\beta-1+
 \frac{\log\!\left((L_{q,\beta}-y)/(L_{q,\beta}[\beta]_q)\right)}{\log q}
 +o(1).
 \label{eq:alpha-inverse-endpoint}
\end{equation}
The inverse diverges only logarithmically as the finite ceiling is approached.

\subsection{Recovering the lower parameter}
\begin{theorem}[\statusproved\ Two lower-parameter branches]
Fix \(\alpha>0\).  The function
\(\beta\mapsto\mathcal G_q(\alpha,\beta)\) is symmetric about
\(\alpha/2\), strictly log-concave on \((0,\alpha)\), equals one at both
endpoints, and has its unique maximum
\[
 M_{q,\alpha}
 =\frac{\Gamma_q(\alpha+1)}{\Gamma_q(\alpha/2+1)^2}
\]
at \(\beta=\alpha/2\).  Every \(y\in(1,M_{q,\alpha})\) therefore has exactly
two inverse branches \(\beta_-<\alpha/2<\beta_+\), related by
\(\beta_+=\alpha-\beta_-\).
\end{theorem}
\begin{proof}
Symmetry is immediate.  Differentiation gives
\[
 \partial_\beta\log\mathcal G_q
 =-\psi_q(\beta+1)+\psi_q(\alpha-\beta+1),
\]
which vanishes only at the center, and
\[
 \partial_\beta^2\log\mathcal G_q
 =-\psi_q'(\beta+1)-\psi_q'(\alpha-\beta+1)<0.
\]
\end{proof}

Near the maximum, with \(x_0=\alpha/2+1\),
\begin{equation}
 \beta_\pm(y)=\frac\alpha2
 \pm\sqrt{\frac{M_{q,\alpha}-y}
 {M_{q,\alpha}\psi_q'(x_0)}}
 +O(M_{q,\alpha}-y).
 \label{eq:beta-branches-max}
\end{equation}
This is another explicit occurrence of the universal square-root inverse law.

\section{Inverting the base in the generalized coefficient}
For nonintegral \(\alpha,\beta\), positivity of polynomial coefficients is no
longer available.  The local derivative is
\[
 \partial_q\log\mathcal G_q
 =\partial_q\log\Gamma_q(\alpha+1)
 -\partial_q\log\Gamma_q(\beta+1)
 -\partial_q\log\Gamma_q(\alpha-\beta+1).
\]
At \(q\to1^-\), the first-order \(q\)-gamma expansion gives
\begin{equation}
 \log\frac{\mathcal G_{\EulerE^{-t}}(\alpha,\beta)}
 {\binom{\alpha}{\beta}}
 =-\frac{\beta(\alpha-\beta)}2t+O(t^2),
 \label{eq:generalized-qbinom-q1}
\end{equation}
where the classical generalized binomial coefficient is interpreted by gamma
functions.  Thus the same effective degree \(D=\beta(\alpha-\beta)\) controls
the leading inverse near \(q=1\).  A global monotonicity theorem for all real
\(\alpha\ge\beta\ge0\) would be valuable; the integer theorem proves it on the
lattice but does not automatically interpolate.

\begin{researchproblem}
Determine the maximal region in the \((\alpha,\beta)\)-plane on which
\(q\mapsto\mathcal G_q(\alpha,\beta)\) is strictly increasing for \(q>0\), or
strictly increasing on \((0,1)\).  Express any boundary in terms of a completely
monotone Lambert kernel.
\end{researchproblem}

\chapter{A positive-coefficient inversion principle}

\section{The general theorem}
\begin{theorem}[\statusproved\ Positive-polynomial inverse principle]
Let
\[
 F(q)=c_0+c_1q+\cdots+c_Dq^D,
 \qquad c_j\ge0,
\]
with at least one \(c_j>0\) for \(j\ge1\).  Then \(F\) is strictly increasing
on \((0,\infty)\), from \(c_0\) to \(+\infty\).  Hence it has a unique positive
inverse for every \(y>c_0\).  The inverse is analytic except at images of
complex critical points and has condition number
\[
 \kappa(q)=\AbsoluteValue{\frac{F(q)}{qF'(q)}}.
\]
\end{theorem}
\begin{proof}
For \(q>0\), \(F'(q)=\sum_{j\ge1}jc_jq^{j-1}>0\).  The remaining claims follow
from the inverse function theorem.
\end{proof}

This elementary result is a powerful inverse theorem because many combinatorial
\(q\)-analogs are polynomials with nonnegative coefficients.

\section{Examples}
\subsection{\(q\)-factorials}
For \(n\ge2\),
\[
 [n]_q!=\prod_{j=1}^n(1+q+\cdots+q^{j-1})
\]
is a nonconstant polynomial with nonnegative coefficients.  It has one
positive-base inverse for every \(y>1\); on \((0,1)\), the range is
\((1,n!)\).

\subsection{\(q\)-multinomial coefficients}
For a composition \(n_1+\cdots+n_r=n\),
\[
 \GaussianBinomial{n}{n_1,\ldots,n_r}{q}
 =\frac{[n]_q!}{[n_1]_q!\cdots[n_r]_q!}
\]
is a polynomial with nonnegative coefficients.  Unless the composition is
trivial, every target greater than one has one positive-base inverse.  Near
\(q=1\), its logarithm is the cumulant-generating function of the multiset
inversion statistic, so mean and variance determine the first inverse terms.

\subsection{Rogers--Szego polynomials}
For fixed \(z>0\),
\[
 H_n(z;q)=\sum_{k=0}^n\GaussianBinomial{n}{k}{q} z^k
\]
has nonnegative coefficients as a polynomial in \(q\).  It therefore has a
unique positive-base inverse for every target above \(H_n(z;0)\), unless
\(n\le1\) makes it constant in \(q\).

\subsection{Partition polynomials and selected \(q\)-Catalan numbers}
Any finite partition generating polynomial
\(\sum_{\lambda\in\mathscr P}q^{w(\lambda)}\), where \(\mathscr P\) is the
specified finite family of partitions, with a nonconstant nonnegative
weight has a unique positive-base inverse.  The same conclusion applies to
those normalizations of \(q\)-Catalan numbers known to be polynomials with
nonnegative coefficients.  Because several inequivalent \(q\)-Catalan
families exist, positivity should be checked for the chosen definition rather
than inferred from the name alone.

\section{Probabilistic inversion}
Normalize \(F\) at \(q=1\) and let \(X\) have probabilities
\(\Pr(X=j)=c_j/F(1)\).  Then
\[
 \log\frac{F(\EulerE^t)}{F(1)}
 =\log\mathbb E\EulerE^{tX}
 =\sum_{r\ge1}\kappa_r\frac{t^r}{r!},
\]
where \(\kappa_r\) are cumulants.  Inverting a target near \(F(1)\) is thus
series reversion of a cumulant expansion.  For palindromic polynomials, all odd
centered cumulants beyond the first vanish, explaining the missing cubic term
in the Gaussian forward expansion and the simple inverse coefficients in
\Cref{eq:qbinom-inverse-one}.

\begin{principle}[Cumulant-to-inverse pipeline]
For any positive-coefficient \(q\)-analogue with a known combinatorial
statistic, compute its cumulants, form the logarithmic target displacement, and
revert the cumulant series.  This separates combinatorics (cumulants) from
analysis (series inversion) and applies uniformly to Gaussian, multinomial,
Mahonian, and partition polynomials.
\end{principle}

\chapter{Inverse \texorpdfstring{\(q\)}{q}-numbers and elementary deformations}

\section{The \texorpdfstring{\(q\)}{q}-number}
For \(q>0\), \(q\ne1\), define
\[
 [x]_q=\frac{1-q^x}{1-q}.
\]

\begin{theorem}[\statusproved\ Exact inverse in \(x\)]
For a fixed branch of the logarithm and an admissible target,
\begin{equation}
 x=\frac{\PrincipalLogarithm(1-(1-q)y)}{\PrincipalLogarithm q}.
 \label{eq:qnumber-inverse-x}
\end{equation}
For real \(0<q<1\), this gives a strictly increasing bijection
\(x\in(0,\infty)\mapsto y\in(0,(1-q)^{-1})\).
\end{theorem}
\begin{proof}
Rearrange to \(q^x=1-(1-q)y\) and take a chosen logarithm.  Real monotonicity
follows from \(q^x\) decreasing in \(x\).
\end{proof}

As a complex inverse, \eqref{eq:qnumber-inverse-x} has the logarithmic lattice
\[
 x_m=\frac{\PrincipalLogarithm(1-(1-q)y)+2\pi\ImaginaryUnit m}{\PrincipalLogarithm q}.
\]

\section{Recovering the base of a \texorpdfstring{\(q\)}{q}-number}
\begin{theorem}[\statusproved\ Real base monotonicity]
Fix \(x>0\), \(x\ne1\).  On \(0<q<1\), the function \([x]_q\) is strictly
increasing from \(1\) to \(x\) if \(x>1\), and strictly decreasing from \(1\)
to \(x\) if \(0<x<1\).  Thus every target strictly between \(1\) and \(x\)
has one base inverse.
\end{theorem}
\begin{proof}
Use the integral average
\[
 [x]_q=\frac1{1-q}\int_q^1xt^{x-1}\Differential t.
\]
If \(x>1\), the integrand is increasing, so deleting the lower part of the
interval raises its average.  If \(0<x<1\), it is decreasing, so the average
falls.  The endpoint limits are immediate.
\end{proof}

For integer \(x=n\), the inverse equation is the polynomial
\[
 1+q+\cdots+q^{n-1}=y.
\]
The positive branch is unique, while the complex inverse has \(n-1\) sheets.

\section{Other elementary \texorpdfstring{\(q\)}{q}-deformations}
The same workflow applies to symmetric \(q\)-numbers
\((q^x-q^{-x})/(q-q^{-1})\), Tsallis-type logarithms, and deformed powers, but
they have different branch symmetries and should not be conflated with the
product-theoretic inverses developed below.  In every case the correct starting
point is the parameter derivative and its zeros, followed by a declared
normalization at \(q=1\) or another regular base.

\part{Inverse \texorpdfstring{\(q\)}{q}-special functions}

\chapter{The inverse \texorpdfstring{\(q\)}{q}-gamma and \texorpdfstring{\(q\)}{q}-polygamma functions}

\section{Two real argument branches of \texorpdfstring{\(\Gamma_q\)}{Gamma-q}}
For \(0<q<1\),
\[
 \Gamma_q(x)=(1-q)^{1-x}\frac{(q;q)_\infty}{(q^x;q)_\infty},
 \qquad x>0.
\]
Its logarithmic derivative is
\begin{equation}
 \psi_q(x)=-\log(1-q)+\log q
 \sum_{j=0}^\infty\frac{q^{x+j}}{1-q^{x+j}},
 \label{eq:qdigamma}
\end{equation}
and
\begin{equation}
 \psi_q'(x)=(\log q)^2
 \sum_{j=0}^\infty\frac{q^{x+j}}{(1-q^{x+j})^2}>0.
 \label{eq:qtrigamma}
\end{equation}

\begin{theorem}[\statusproved\ Unique minimum and two inverse branches]
There is a unique \(x_q^*>0\) satisfying \(\psi_q(x_q^*)=0\).  The function
\(\Gamma_q\) decreases strictly from \(+\infty\) to
\[
 m_q=\Gamma_q(x_q^*)
\]
on \((0,x_q^*)\), then increases strictly to \(+\infty\) on
\((x_q^*,\infty)\).  Thus every \(y>m_q\) has exactly two positive inverse
values
\[
 \Gamma_{q,-}^{-1}(y)<x_q^*<\Gamma_{q,+}^{-1}(y),
\]
while \(y=m_q\) has one double inverse value.
\end{theorem}
\begin{proof}
Equation \eqref{eq:qtrigamma} makes \(\psi_q\) strictly increasing.  As
\(x\downarrow0\), its first summand yields \(\psi_q(x)=-1/x+O(1)\), so the
limit is \(-\infty\).  As \(x\to\infty\), the sum tends to zero and
\(\psi_q(x)\to-\log(1-q)>0\).  Thus there is exactly one zero.  Since
\(\Gamma_q'=\Gamma_q\psi_q\), this is the unique minimum.  The product formula
shows \(\Gamma_q(x)\to\infty\) at both endpoints.
\end{proof}

\begin{figure}[htbp]
\centering
\includegraphics[width=0.88\textwidth]{figures/qgamma_inverse_branches.pdf}
\caption{The unique minimum and the two positive inverse branches of
\(\Gamma_q\).}
\label{fig:qgamma-branches}
\end{figure}

\section{Puiseux expansion at the minimum}
Let
\[
 \tau_q=\psi_q'(x_q^*)>0,
 \qquad
 \sigma_q=\psi_q''(x_q^*)<0.
\]
At the minimum,
\[
 \Gamma_q''(x_q^*)=m_q\tau_q,
 \qquad
 \Gamma_q'''(x_q^*)=m_q\sigma_q.
\]
Hence
\begin{equation}
 \Gamma_{q,\pm}^{-1}(y)
 =x_q^*\pm\sqrt{\frac{2(y-m_q)}{m_q\tau_q}}
 -\frac{\sigma_q}{3m_q\tau_q^2}(y-m_q)
 +O((y-m_q)^{3/2}).
 \label{eq:qgamma-minimum-inverse}
\end{equation}
This expansion is ideal for initializing both branches when the target is close
to the minimum, where direct Newton iteration is ill-conditioned.

\section{Endpoint asymptotics of the two branches}
Put
\[
 \lambda=-\log q,
 \qquad
 L=-\log(1-q),
 \qquad
 C_q=(1-q)(q;q)_\infty.
\]
As \(x\downarrow0\),
\[
 \Gamma_q(x)=\frac{1-q}{\lambda x}(1+O(x)),
\]
so
\begin{equation}
 \Gamma_{q,-}^{-1}(y)
 =\frac{1-q}{\lambda y}\bigl(1+O(y^{-1})\bigr),
 \qquad y\to\infty.
 \label{eq:qgamma-lower-large}
\end{equation}
As \(x\to\infty\),
\[
 \Gamma_q(x)=C_q\EulerE^{Lx}
 \left(1+\frac{q^x}{1-q}+O(q^{2x})\right),
\]
which yields
\begin{equation}
\begin{split}
 \Gamma_{q,+}^{-1}(y)
 ={}&\frac{\log(y/C_q)}{L}\\
 &-\frac{1}{L(1-q)}
 \left(\frac{C_q}{y}\right)^{\lambda/L}
 +O\!\left(y^{-2\lambda/L}+y^{-\lambda/L-1}\right).
\end{split}
\label{eq:qgamma-upper-large}
\end{equation}
The two branches therefore have radically different large-target scales:
reciprocal on the left and logarithmic on the right.

\section{Recovering the base of \texorpdfstring{\(\Gamma_q(x)\)}{Gamma-q(x)}}
The dependence on \(q\) is governed by a classical monotonicity theorem
recorded in the DLMF \cite{DLMFqGamma}.

\begin{theorem}[Known monotonicity; \statusproved]
Fix \(x>0\).
\begin{enumerate}[label=(\roman*)]
\item If \(0<x<1\) or \(x>2\), then \(q\mapsto\Gamma_q(x)\) is strictly
increasing on \((0,1)\), from \(1\) to \(\Gamma(x)\).
\item If \(1<x<2\), it is strictly decreasing from \(1\) to \(\Gamma(x)\).
\item At \(x=1\) and \(x=2\), it is identically one.
\end{enumerate}
Consequently the base inverse is unique whenever the target lies strictly
between \(1\) and \(\Gamma(x)\), with the order of the endpoints determined by
the stated sector.
\end{theorem}

At \(q\to0^+\),
\begin{equation}
 \log\Gamma_q(x)=q^x+(x-2)q+o(q^x+q),
 \label{eq:qgamma-smallq}
\end{equation}
away from the exact exceptional points \(x=1,2\).  Therefore
\begin{align*}
 q&\sim(y-1)^{1/x},&&0<x<1,\\
 q&\sim\frac{1-y}{2-x},&&1<x<2,\\
 q&\sim\frac{y-1}{x-2},&&x>2.
\end{align*}
At the classical endpoint, with \(q=\EulerE^{-t}\),
\begin{equation}
 \log\Gamma_q(x)
 =\log\Gamma(x)-\frac{(x-1)(x-2)}4t+O(t^2).
 \label{eq:qgamma-q1-first}
\end{equation}
Thus, if \(y\to\Gamma(x)\) within the appropriate monotonic range,
\begin{equation}
 -\log q
 =\frac{4\log(\Gamma(x)/y)}{(x-1)(x-2)}
 +O\bigl(\log(\Gamma(x)/y)^2\bigr).
 \label{eq:qgamma-base-inverse-q1}
\end{equation}
The quotient is positive in all three nonconstant sectors.

\section{The minimizer as a function of the base}
The numerical data in the archive show
\[
 x_{0.05}^*=1.37998346149\ldots,
 \qquad
 x_{0.95}^*=1.46060675866\ldots,
\]
while the classical gamma minimizer is
\(x^*=1.461632144968\ldots\).

For \(q\to0^+\), let \(\lambda=-\log q\).  Balancing the leading terms in
\(\psi_q(x)=0\), namely \(q-\lambda q^x=0\), predicts
\begin{equation}
 x_q^*=1+\frac{\log\lambda}{\lambda}+o(\lambda^{-1}).
 \label{eq:qgamma-minimizer-smallq}
\end{equation}
For \(q\to1^-\), local uniform convergence of \(\psi_q\) to the classical
\(\psi\) implies \(x_q^*\to x^*\).

\begin{conjecture}[\statusconjectural\ Monotone minimizer]
The function \(q\mapsto x_q^*\) is strictly increasing on \((0,1)\), and the
minimum value \(m_q\) is strictly decreasing from \(1\) to
\(\Gamma(x^*)\).
\end{conjecture}
The data in \texttt{data/qgamma\_minima.csv} support both assertions.  A proof
should follow from a sign analysis of
\[
 \frac{\Differential x_q^*}{\Differential q}
 =-\frac{\partial_q\psi_q(x_q^*)}{\psi_q'(x_q^*)},
\]
so the problem reduces to the sign of one explicit Lambert-series derivative.

\section{Inverse \texorpdfstring{\(q\)}{q}-digamma and
\texorpdfstring{\(q\)}{q}-polygamma functions}
Expanding the geometric denominator gives
\begin{equation}
 \psi_q(x)=-\log(1-q)+\log q
 \sum_{\ell=1}^\infty\frac{q^{\ell x}}{1-q^\ell},
 \label{eq:qdigamma-lambert}
\end{equation}
and, for \(m\ge1\),
\begin{equation}
 \psi_q^{(m)}(x)=(\log q)^{m+1}
 \sum_{\ell=1}^\infty\frac{\ell^m q^{\ell x}}{1-q^\ell}.
 \label{eq:qpolygamma-lambert}
\end{equation}

\begin{theorem}[\statusproved\ Global inverse polygamma branches]
For fixed \(0<q<1\):
\begin{enumerate}[label=(\roman*)]
\item \(\psi_q\) is a strictly increasing bijection from \((0,\infty)\) onto
\(( -\infty,-\log(1-q))\).  It has one real inverse for every target in this
range.
\item For every \(m\ge1\), the positive function
\[
 g_{q,m}(x)=(-1)^{m+1}\psi_q^{(m)}(x)
\]
is strictly decreasing from \(+\infty\) to \(0\).  Hence it has one real
inverse for every positive target.
\end{enumerate}
\end{theorem}
\begin{proof}
The first assertion follows from \(\psi_q'>0\) and its endpoint limits.  Formula
\eqref{eq:qpolygamma-lambert} shows \(g_{q,m}>0\); differentiating gives
\(g_{q,m}'=-g_{q,m+1}<0\).  The endpoint limits follow from the first terms and
comparison with the corresponding powers of \(x^{-1}\).
\end{proof}

This family provides unusually clean examples of globally invertible
transcendental \(q\)-functions.

\chapter{The inverse \texorpdfstring{\(q\)}{q}-beta function}

\section{Coordinate inverses}
For \(a,b>0\),
\[
 B_q(a,b)=\frac{\Gamma_q(a)\Gamma_q(b)}{\Gamma_q(a+b)}.
\]

\begin{theorem}[\statusproved\ Global inverse in either shape parameter]
Fix \(0<q<1\) and \(b>0\).  The function \(a\mapsto B_q(a,b)\) is strictly
decreasing and log-convex on \((0,\infty)\), from \(+\infty\) to
\begin{equation}
 \ell_{q,b}=(1-q)^b\Gamma_q(b).
 \label{eq:qbeta-floor}
\end{equation}
Therefore every \(y>\ell_{q,b}\) has exactly one inverse \(a\).  The same
statement holds with \(a\) and \(b\) interchanged.
\end{theorem}
\begin{proof}
Differentiation gives
\[
 \partial_a\log B_q(a,b)=\psi_q(a)-\psi_q(a+b)<0,
\]
and
\[
 \partial_a^2\log B_q(a,b)=\psi_q'(a)-\psi_q'(a+b)>0
\]
because \(\psi_q\) increases and \(\psi_q'\) decreases.  The limit at zero is
infinite because \(\Gamma_q(a)\) has a simple pole.  The product formula gives
\[
 B_q(a,b)=
 (1-q)^b\Gamma_q(b)\frac{(q^{a+b};q)_\infty}{(q^a;q)_\infty}
 =\frac{\ell_{q,b}}{(q^a;q)_b},
\]
which tends to \(\ell_{q,b}\) as \(a\to\infty\).
\end{proof}

The Hessian of \(\log B_q\) is
\[
 \begin{pmatrix}
 \psi_q'(a)-\psi_q'(a+b)&-\psi_q'(a+b)\\
 -\psi_q'(a+b)&\psi_q'(b)-\psi_q'(a+b)
 \end{pmatrix}.
\]
Coordinate convexity is immediate; joint convexity requires positivity of the
determinant and is a separate question.

\section{Endpoint inverse asymptotics}
As \(a\downarrow0\),
\[
 B_q(a,b)=\frac{1-q}{(-\log q)a}(1+O(a)),
\]
so
\begin{equation}
 a(y)=\frac{1-q}{(-\log q)y}\bigl(1+O(y^{-1})\bigr),
 \qquad y\to\infty.
 \label{eq:qbeta-inverse-large}
\end{equation}
At the other endpoint,
\[
 B_q(a,b)=\ell_{q,b}
 \bigl(1+[b]_qq^a+O(q^{2a})\bigr),
\]
and therefore
\begin{equation}
 a(y)=
 \frac{\log\!\left((y-\ell_{q,b})/(\ell_{q,b}[b]_q)\right)}{\log q}
 +o(1),
 \qquad y\downarrow\ell_{q,b}.
 \label{eq:qbeta-inverse-floor}
\end{equation}
Again, approach to a finite floor produces logarithmic divergence of the
inverse parameter.

\section{Inversion with respect to \texorpdfstring{\(q\)}{q}}
The product representation
\[
 B_q(a,b)=(1-q)
 \frac{(q;q)_\infty(q^{a+b};q)_\infty}
 {(q^a;q)_\infty(q^b;q)_\infty}
\]
gives the exact kernel
\begin{align}
 q\partial_q\log B_q(a,b)
={}&-\frac{q}{1-q}
 -\sum_{j=1}^\infty\frac{jq^j}{1-q^j}
 -\sum_{j=0}^\infty\frac{(a+b+j)q^{a+b+j}}{1-q^{a+b+j}}\notag\\
&+\sum_{j=0}^\infty\frac{(a+j)q^{a+j}}{1-q^{a+j}}
 +\sum_{j=0}^\infty\frac{(b+j)q^{b+j}}{1-q^{b+j}}.
 \label{eq:qbeta-base-kernel}
\end{align}
Unlike the coordinate derivatives, its sign is not transparent for arbitrary
real \(a,b\).  The endpoint values are
\[
 \lim_{q\to0^+}B_q(a,b)=1,
 \qquad
 \lim_{q\to1^-}B_q(a,b)=B(a,b).
\]
Depending on whether the classical beta value is above or below one, an
increasing or decreasing trend is possible; internal extrema are not excluded
without a sign theorem for \eqref{eq:qbeta-base-kernel}.

For positive integer parameters, however, the answer is complete.

\begin{theorem}[\statusproved\ Integer-parameter base inverse]
Let \(m,n\in\NaturalNumbers\), not both equal to one.  Then
\begin{equation}
 B_q(m,n)=
 \frac1{[n]_q\GaussianBinomial{m+n-1}{m-1}{q}}.
 \label{eq:qbeta-integer-reciprocal}
\end{equation}
It is strictly decreasing for \(q>0\), from \(1\) to \(0\).  Hence every
\(y\in(0,1)\) has exactly one positive-base inverse.  Restricted to
\(0<q<1\), the target range is \((B(m,n),1)\).
\end{theorem}
\begin{proof}
Use \(\Gamma_q(r)=[r-1]_q!\) and simplify to
\eqref{eq:qbeta-integer-reciprocal}.  The denominator is a nonconstant product
of positive-coefficient polynomials and is strictly increasing for \(q>0\).
\end{proof}

\begin{conjecture}[\statusconjectural\ Real-parameter beta sign regions]
There is a finite collection of explicit curves in the positive
\((a,b)\)-quadrant separating regions on which
\(q\mapsto B_q(a,b)\) is monotone on \((0,1)\).  In particular, it is expected
to be decreasing throughout \(a\ge1,b\ge1\), extending the integer theorem.
\end{conjecture}
A proof or counterexample should begin with complete-monotonicity transforms of
\eqref{eq:qbeta-base-kernel}, not with direct numerical differentiation.

\chapter{Euler \texorpdfstring{\(q\)}{q}-exponentials and product logarithms}

\section{Exact inverse reductions}
Use the unscaled Euler conventions
\[
 e_q(z)=\frac1{(z;q)_\infty},
 \qquad
 E_q(z)=(-z;q)_\infty,
 \qquad 0<q<1.
\]
(The common scaled conventions replace \(z\) by \((1-q)z\).)

\begin{theorem}[\statusproved\ Exact real inverses]
On their principal real branches,
\begin{align}
 e_q^{-1}(y)&=\Aq(1/y),&&y>0,\label{eq:eq-inverse}\\
 E_q^{-1}(y)&=-\Aq(y),&&y>0.\label{eq:Eq-inverse}
\end{align}
For the scaled functions
\(\widetilde e_q(z)=1/((1-q)z;q)_\infty\) and
\(\widetilde E_q(z)=(-(1-q)z;q)_\infty\), divide the right sides by
\(1-q\).
\end{theorem}
\begin{proof}
Equations \eqref{eq:eq-inverse}--\eqref{eq:Eq-inverse} are immediate from the
definitions and the uniqueness of the principal product inverse.
\end{proof}

This suggests product-theoretic \(q\)-logarithms
\[
 \log_q^{(e)}(y)=\Aq(1/y),
 \qquad
 \log_q^{(E)}(y)=-\Aq(y).
\]
They are distinct from Tsallis and Jackson logarithms.  Their analytic theory is
inherited exactly from \(\Aq\), including the Bernstein conjecture, branch
points, and endpoint expansions.

\section{Asymptotic consequences}
From \eqref{eq:Aq-near-zero}, as \(y\to\infty\),
\[
 e_q^{-1}(y)=1-\frac1{(q;q)_\infty y}
 +\frac{L_q}{(q;q)_\infty^2y^2}+O(y^{-3}).
\]
From \eqref{eq:Aq-large-leading}, as \(y\to\infty\),
\[
 \log E_q^{-1}(y)
 =\sqrt{2(-\log q)\log y}+\frac12\log q
 +O((\log y)^{-1/2}).
\]
Thus the two Euler exponentials have inverse scales approaching a finite
singularity and growing stretched-exponentially, respectively.

\section{Inverse semigroup questions}
The ordinary logarithm linearizes multiplication.  Neither product logarithm
above does so exactly.  One may instead define a deformed operation
\[
 x\oplus_q y=e_q^{-1}(e_q(x)e_q(y)),
\]
on a branch where the product remains admissible.  Associativity follows from
ordinary multiplication, while the explicit operation is expressed through
\(\Aq\).  Its Taylor coefficients, singularities, and classical limit form a
concrete inverse-function approach to \(q\)-addition.

\begin{researchproblem}
Classify the real and complex Lie-semigroup structure induced by
\(\oplus_q\), determine whether its infinitesimal generator has a simple
Lambert-series representation, and compare it with Jackson's and Tsallis'
standard deformed additions.
\end{researchproblem}

\chapter{Basic hypergeometric, theta, and orthogonal-polynomial inverses}

\section{Local inverse of a basic hypergeometric series}
In a standard convention,
\[
 {}_r\phi_s\!\left(
 \begin{matrix}a_1,\ldots,a_r\\b_1,\ldots,b_s\end{matrix};q,z\right)
 =\sum_{m\ge0}c_mz^m,
\]
where
\[
 c_m=\frac{(a_1,\ldots,a_r;q)_m}
 {(q,b_1,\ldots,b_s;q)_m}
 \left((-1)^mq^{\binom m2}\right)^{1+s-r}.
\]
Assume \(c_1\ne0\).  Since the value at \(z=0\) is one, the local inverse in
the argument is
\begin{equation}
 z=\frac{y-1}{c_1}
 -\frac{c_2}{c_1^3}(y-1)^2
 +\left(\frac{2c_2^2}{c_1^5}-\frac{c_3}{c_1^4}\right)(y-1)^3
 +O((y-1)^4).
 \label{eq:basic-hypergeometric-inverse}
\end{equation}
If \(c_1=0\) but \(c_r\ne0\) is the first nonzero coefficient, the inverse has
\(r\)-th-root ramification at \(y=1\).

Parameter inverses are obtained from the same implicit formula
\[
 \frac{\partial p}{\partial y}=\frac1{F_p},
 \qquad
 \frac{\partial p}{\partial a_i}=-\frac{F_{a_i}}{F_p},
\]
where every parameter derivative of a coefficient reduces to finite or
infinite \(q\)-Pochhammer kernels.  Singular divisors occur when denominator
factors vanish, and branch points occur when the chosen parameter derivative
vanishes.

\section{Theta-function inversion}
Define
\[
 \theta(z;q)=(z;q)_\infty(q/z;q)_\infty.
\]
It satisfies
\[
 \theta(qz;q)=-z^{-1}\theta(z;q)
\]
and has zeros at \(z=q^m\), \(m\in\IntegerNumbers\).  Its logarithmic argument derivative is
\begin{equation}
 z\partial_z\log\theta(z;q)
 =-\sum_{j\ge0}\frac{zq^j}{1-zq^j}
 +\sum_{j\ge0}\frac{q^{j+1}/z}{1-q^{j+1}/z}.
 \label{eq:theta-critical-kernel}
\end{equation}
Thus inverse theta branch points are the images of zeros of the balanced
Lambert kernel \eqref{eq:theta-critical-kernel}.  Quasi-periodicity transports
local data between annuli, while the two product factors make the inverse
surface doubly infinite.

The phase function in \eqref{eq:logperiodic-product} is itself a theta-product
factor.  Large-target inverse \(q\)-Pochhammer asymptotics and theta inversion
are therefore two aspects of the same multiplicative lattice geometry.

\section{\texorpdfstring{\(q\)}{q}-orthogonal polynomials}
For fixed degree and parameters, an orthogonal polynomial
\(p_n(x;q)\) has an algebraic inverse in \(x\).  Its real branches are divided
by its critical points, which interlace the zeros under the usual positive
orthogonality hypotheses.  The branch-value polynomial is
\[
 \Disc_x(p_n(x;q)-y).
\]
Inversion in a shape parameter or in \(q\) is usually transcendental even when
inversion in \(x\) is algebraic.  Parameter derivatives can be obtained from
three-term recurrences by differentiating the recurrence and propagating both
\(p_n\) and its sensitivity.

\begin{principle}[Recurrence-assisted inverse continuation]
For a family satisfying
\[
 p_{n+1}(x)=(A_nx+B_n)p_n(x)-C_np_{n-1}(x),
\]
propagate \(p_n\), \(\partial_xp_n\), and all needed parameter derivatives
through differentiated recurrences.  Use these quantities in a branch-aware
Newton corrector.  This avoids unstable coefficient expansion at large degree.
\end{principle}

\section{\texorpdfstring{\(q\)}{q}-Bessel and
\texorpdfstring{\(q\)}{q}-Airy functions}
In standard parameter regimes, Jackson \(q\)-Bessel and \(q\)-Airy functions
may exhibit infinitely many zeros and critical values.  When that structure is
present, their argument inverses are naturally indexed by zero intervals on
the real axis and by continuation paths in the complex plane.  As with the
ordinary Bessel inverse, a single principal branch captures only a small part
of the global surface.  Uniform asymptotics for large zero index can be
inverted to initialize remote branches; near coalescing turning points,
Airy-type uniformization replaces ordinary Taylor reversion.

\begin{researchproblem}
Construct a DLMF-style branch atlas for inverse Jackson \(q\)-Bessel and
\(q\)-Airy functions, including zero-index labels, critical values,
\(q\to1\) transition formulas, and certified continuation algorithms.  The
product-kernel methods of this report supply the parameter derivatives needed
for such an atlas.
\end{researchproblem}

\part{Computation, certification, and research frontiers}

\chapter{Branch-aware numerical inversion}

\section{Why unconstrained Newton iteration is insufficient}
Newton's update for \(F(p)=y\) is
\[
 p_{m+1}=p_m-\frac{F(p_m)-y}{F'(p_m)}.
\]
Near a simple root it converges quadratically, but globally it can cross poles,
jump over critical points, or converge to a different sheet.  A reliable
inverse routine must carry a branch certificate: a monotone interval on the
real axis, or a continuation path and sheet label in the complex plane.

\section{Bracketed Newton on a real branch}
\begin{algorithm}[\statusproved\ Safeguarded branch inversion]
Suppose \(F\in C^2([L,R])\) is strictly monotone,
\(\lvert F'(p)\rvert\ge m>0\) throughout the interval, and \(y\) lies between
\(F(L)\) and \(F(R)\).
\begin{enumerate}[label=\arabic*.]
\item Maintain a bracket \([L_m,R_m]\) known to contain the unique inverse.
\item Propose the Newton point
\[
 N_m=p_m-\frac{F(p_m)-y}{F'(p_m)}.
\]
\item Accept \(N_m\) only if it lies in the middle half
\([L_m+(R_m-L_m)/4,R_m-(R_m-L_m)/4]\); otherwise replace it by the midpoint.
\item Update the bracket using monotonicity, set \(p_{m+1}=N_m\), and stop when
both the residual and bracket width meet tolerance.
\end{enumerate}
The bracket never loses the solution and its width contracts by at least a
factor \(3/4\) per step.  Once unsafeguarded Newton points remain accepted in
the local basin, the iterates converge quadratically.
\end{algorithm}

This algorithm is used in the accompanying code for the principal
\(q\)-Pochhammer inverse, finite and infinite critical points, Gaussian base
inverses, and the two \(q\)-gamma branches.

\section{Stable product evaluation}
For positive real factors, evaluate
\[
 \log(a;q)_\infty=\sum_{j\ge0}\log(1-aq^j)
\]
rather than multiplying directly.  Use \texttt{log1p} for small \(aq^j\), and
truncate only after applying the tail bound of \Cref{thm:product-tail}.
For signed real intervals, separate the sign of the finite number of negative
factors and accumulate the logarithm of the absolute value.

For derivatives, logarithmic kernels are usually more stable than direct
finite differences.  At a principal inverse point \(a=\Aq(y)\), for example,
\[
 \Aq'(y)=-\frac1{yS_1(a)}
\]
requires only a positive Lambert sum.

\section{Complex predictor--corrector continuation}
Suppose a branch value \(p(y_0)\) is known and a path
\(y(s)\), \(0\le s\le1\), avoids critical values.  Differentiate
\(F(p(s))=y(s)\):
\[
 p'(s)=\frac{y'(s)}{F'(p(s))}.
\]
A first-order predictor is
\[
 p_{\mathrm{pred}}=p_m+\frac{\Delta y}{F'(p_m)}.
\]
Correct it with Newton's method at the new target.  Step size should be reduced
when \(\AbsoluteValue{F'}\) becomes small or the corrector takes too many iterations.
Near a known branch point, switch to the local square-root coordinate
\(w=\sqrt{y-y_c}\), in which the sheets become analytic.

\begin{algorithm}[Complex inverse continuation]
\begin{enumerate}[label=\arabic*.]
\item Choose a base point, a branch root, and a target path avoiding the
critical-value set.
\item Predict using the implicit differential equation.
\item Correct with Newton iteration and verify the residual.
\item Compare the corrected root with all nearby known roots to detect an
unintended sheet jump.
\item Adapt the step according to \(\AbsoluteValue{F'}\), Newton contraction, and distance
to known critical values.
\item Record the permutation after closed loops to compute monodromy.
\end{enumerate}
\end{algorithm}

\section{Locating branch points}
For a one-variable direct map, solve
\[
 F_p(p;\lambda)=0,
 \qquad
 y=F(p;\lambda).
\]
For finite algebraic families, eliminate \(p\) using a resultant or
discriminant.  To find parameter values where branch values collide, compute a
second discriminant in \(y\).  This is exactly how
\eqref{eq:H5} was found.

For transcendental products, bracket the real critical point using poles of the
logarithmic derivative.  In the complex plane, use argument-principle counts to
ensure that all critical points in a box or disk have been found:
\[
 N=\frac1{2\pi\ImaginaryUnit}\int_\gamma\frac{F_{pp}(p)}{F_p(p)}\Differential p.
\]
Subdivide until each region contains one root, then apply interval Newton or
ball arithmetic.

\section{Formal and interval certification}
The following hierarchy separates levels of trust.
\begin{enumerate}[label=(\roman*)]
\item \emph{Floating-point evidence}: useful for discovery, not proof.
\item \emph{Arbitrary-precision residual}: rules out ordinary roundoff but not a
systematic branch or truncation error.
\item \emph{Interval or ball proof}: encloses products, derivatives, and roots
with directed rounding.
\item \emph{Exact algebra}: proves polynomial identities and isolating intervals
using rational arithmetic and Sturm sequences.
\item \emph{Proof-assistant formalization}: verifies definitions, inequalities,
and algorithms inside a foundational kernel.
\end{enumerate}
The secondary-discriminant factorizations in the archive are exact algebra.
The decimal identification of colliding critical intervals is high-precision
numerics and is explicitly labelled as such.

\section{Complexity considerations}
A truncated infinite product with \(N\) factors costs \(O(N)\).  As \(q\to1\),
the tail bound in \Cref{thm:product-tail} gives, for fixed nonzero \(a\),
\[
 N=\Theta\!\left(\frac{\log(1/\varepsilon)+\log(1/(1-q))}{1-q}\right)
\]
in the naive fixed-tolerance regime.  Polylogarithmic asymptotics or modular
transformations therefore become essential.
Finite Gaussian coefficients should be evaluated through logarithmic products
or recurrences rather than expanded coefficients when \(n\) is large.

For many targets at fixed parameters, precompute critical points and build a
branch atlas.  Chebyshev or rational approximants on each monotone interval can
then reduce repeated inversion to a few Newton corrections.

\chapter{Reproducible experiments}

\section{Archive layout}
The companion program is \path{numerical_experiments.py}.  It uses only
\texttt{mpmath}, \texttt{sympy}, \texttt{numpy}, and \texttt{matplotlib},
and it creates:
\begin{itemize}
\item exact secondary discriminants for \(3\le n\le6\);
\item finite critical points and values for several \(n,q\);
\item infinite critical-point asymptotic ratios;
\item Gaussian inverse errors near \(q=1\);
\item \(q\)-gamma minimizers;
\item inverse-product derivatives through order eight;
\item all figures included in this report.
\end{itemize}
No external network calls are made.

\section{Experiment A: the first critical-value collision}
The symbolic calculation forms
\[
 P_n(a;q)=\prod_{j=0}^{n-1}(1-aq^j),
\quad
 \Delta_n(y;q)=\Disc_a(P_n-y),
\quad
 \mathfrak M_n(q)=\Disc_y\Delta_n.
\]
Factoring over \(\IntegerNumbers[q]\) produces \eqref{eq:M3}--\eqref{eq:H5}.  Bisection on
the exact factor \(H_5\) gives the quoted \(q_0\).  Independent bracketed
solutions of the logarithmic critical equation locate \(c_{5,0}\) and
\(c_{5,2}\).

A useful lesson is methodological: a plausible numerical conjecture that
critical-value magnitudes increase with the interval index is false near
\(q=1\).  Exact elimination revealed the first crossing and replaced the false
conjecture with a precise Maxwell-locus theory.

\section{Experiment B: remote critical asymptotics}
For each \(q\in\{0.35,0.55,0.75\}\) and \(1\le k\le14\), the program solves
\[
 \sum_{j\ge0}\frac{q^j}{1-aq^j}=0
\]
on \((q^{-k},q^{-(k+1)})\).  It compares the result with
\eqref{eq:ck-asymptotic} and \eqref{eq:vk-asymptotic}.  At \(q=0.35\), for
example, the critical-point ratio ``exact/asymptotic'' improves from about
\(1.0388\) at \(k=2\) to \(1.00056\) at \(k=14\), while the critical-value
ratio improves from about \(0.954\) to \(0.996\).

The convergence slows as \(q\to1\), which is consistent with loss of uniformity
when the product tail length grows.  A double asymptotic in \(k\) and
\(1-q\) is therefore a natural next step.

\section{Experiment C: inverse Bernoulli expansion}
For \(n=12\), \(k=5\), and targets
\(y=\binom{12}{5}\EulerE^{-s}\), the exact positive base is found by bisection.  At
\(s=10^{-4}\), the cubic approximation in
\eqref{eq:qbinom-inverse-one} has relative error below \(10^{-15}\) in
\(t=-\log q\), whereas the quadratic error is about
\(7.7\times10^{-11}\).  The log-log slopes in
\Cref{fig:qbinom-inverse-accuracy} exhibit the predicted orders.

\section{Experiment D: the \texorpdfstring{\(q\)}{q}-gamma minimum}
The root of \(\psi_q(x)=0\) is found by safeguarded Newton iteration using the
positive trigamma derivative.  Selected values are
\begin{center}
\begin{tabular}{@{}ccc@{}}
\toprule
\(q\)&\(x_q^*\)&\(m_q\)\\
\midrule
0.05&1.37998346149&0.98270463141\\
0.20&1.42196580342&0.95376342969\\
0.60&1.45066979499&0.91169489999\\
0.95&1.46060675866&0.88840573828\\
\bottomrule
\end{tabular}
\end{center}
The trend supports the monotone-minimizer conjecture and approaches the
classical gamma minimum.

\section{Experiment E: higher derivatives of the principal inverse}
At a chosen direct argument \(a<1\), write
\[
 (a+h;q)_\infty=(a;q)_\infty
 \prod_{j\ge0}(1-hw_j),
 \qquad
 w_j=\frac{q^j}{1-aq^j}.
\]
Newton identities recover the Taylor coefficients from the power sums
\(\sum_jw_j^r\), avoiding unstable numerical differentiation.  Reverting this
series through degree eight tests the signs of inverse derivatives.  All
archived cases satisfy
\[
 (-1)^r\Aq^{(r)}(y)>0,
 \qquad 1\le r\le8,
\]
which is equivalent to the first eight Bernstein sign conditions for
\(1-\Aq\).

\section{Limitations of the experiments}
The calculations do not prove the Bernstein conjecture, monotonicity of the
\(q\)-gamma minimizer, or the proposed real-parameter beta sign regions.  The
code uses high precision and analytic bracketing, but not directed-rounding
interval arithmetic.  Its purpose is reproducible discovery and stress testing
of formulas.  Every conjectural statement remains labelled.

\chapter{Research program and conjectures}\label{chap:research}

\section{Finite-product Maxwell geometry}
\begin{conjecture}[\statusconjectural\ Structured critical-value crossings]
For fixed \(n\), the positive real zeros of the reduced secondary discriminant
\(\mathfrak M_n(q)\), after removing powers of \(q\), \(q-1\), and cyclotomic
factors, are simple at the level of the collision parameter and correspond to
pairwise equalities of same-sign real critical values.  The number and ordering
of these roots admit a finite combinatorial description in terms of interval
indices.
\end{conjecture}

Questions include:
\begin{enumerate}[label=(\alph*)]
\item Does the first collision always involve \(I_0\) and \(I_2\) for
\(n\ge5\)?
\item How many positive Maxwell parameters occur as \(n\to\infty\)?
\item Do reciprocal collision factors always reduce to polynomials in
\(q+q^{-1}\)?
\item What is the limiting distribution of their roots near \(q=1\)?
\end{enumerate}
Exact resultant computations for moderate \(n\), followed by asymptotic
analysis of critical intervals, form a realistic attack.

\section{A full expansion of remote infinite-product extrema}
\begin{conjecture}[\statusconjectural\ Complete \(1/k\) expansion]
For fixed \(0<q<1\), the critical-point scale in
\eqref{eq:scaled-critical-product} admits expansions
\begin{align*}
 \varepsilon_k
 &\sim\frac1{k+2}+\sum_{r\ge2}\frac{E_r(q)}{(k+2)^r},\\
 c_k
 &\sim q^{-(k+1)}\left(1-\sum_{r\ge1}\frac{C_r(q)}{(k+2)^r}\right),\\
 \frac{\AbsoluteValue{(c_k;q)_\infty}}
 {(q;q)_\infty^2q^{-(k+1)(k+2)/2}}
 &\sim\frac{\EulerE^{-1}}{k+2}
 \left(1+\sum_{r\ge1}\frac{V_r(q)}{(k+2)^r}\right),
\end{align*}
where the coefficients are explicit polynomials in Lambert series such as
\(\sum_{m\ge1}q^m/(1-q^m)^r\).
\end{conjecture}
The proof should expand \(\partial_\varepsilon\log G_m\) uniformly and apply
implicit series reversion.  Exponentially small \(q^m\) terms can then be added
as a transseries correction.

\section{Double scaling and Lambert \texorpdfstring{\(W\)}{W}}
Fixed-\(a\) inversion near \(q=1\) is controlled by
\(\Polylogarithm_2(a)/(-\log q)\).  When \(a\to1\) simultaneously, the dilogarithm and
\(\log(1-a)\) terms compete.  Scaling laws such as
\[
 1-a\asymp -\log q,
 \qquad
 -\log y\asymp\frac1{-\log q},
\]
lead to equations in which the unknown appears both linearly and inside a
logarithm or exponential.  Lambert \(W\), including its \(W_{-1}\) branch, is
then expected to provide uniform transition coordinates.

\begin{researchproblem}
Derive uniform inverse formulas for \((a;q)_\infty=y\) in all joint regimes
\((a,q,y)\to(1,1,0)\).  Identify the transition surfaces between algebraic,
logarithmic, dilogarithmic, and Lambert-\(W\) scales, and prove uniform error
bounds.
\end{researchproblem}

\section{Complete Bernstein and Pick geometry}
A proof that \(1-\Aq\) is complete Bernstein would imply a representation
\[
 1-\Aq(y)=c_0+c_1y+
 \int_0^\infty\frac{y}{y+t}\,\mu_q(\Differential t)
\]
for a positive measure \(\mu_q\).  The singular support should encode the
negative or complex critical values of the direct product.  The challenge is
that inverse-function singularities are not simply the original zero set.

\begin{researchproblem}
Construct, or rule out, a Stieltjes measure for \((1-\Aq(y))/y\).  Relate its
moments to the inverse coefficients at \(y=0\) or \(y=1\), and determine
whether finite products provide monotone approximants in the complete
Bernstein cone.
\end{researchproblem}

\section{Galois and monodromy theory}
For finite \(n\), consider
\[
 P_n(a;q)-y\in\RationalNumbers(q,y)[a].
\]
Away from critical-value collisions, the analytic monodromy is \(S_n\), which
strongly suggests the generic algebraic Galois group over \(\RationalNumbers(q,y)\) is also
\(S_n\).

\begin{conjecture}[\statusconjectural\ Generic Galois group]
For every \(n\ge2\), the polynomial \((a;q)_n-y\) has Galois group \(S_n\) over
\(\RationalNumbers(q,y)\).  Specializations with smaller group are contained in an explicit
thin set controlled by discriminants and resolvents.
\end{conjecture}
A proof can combine irreducibility over the rational function field,
specialization at a convenient rational base, and cycle information from local
ramification.

\section{Inverse-coefficient positivity}
The local inverse of a positive-coefficient polynomial need not itself have
positive coefficients after arbitrary normalization.  Yet the Gaussian
inverse around \(q=0\) exhibits structured integer combinations of partition
numbers, and the principal product inverse exhibits persistent alternating
derivative signs.

\begin{researchproblem}
Find combinatorial interpretations for the coefficients of
\(\mathcal Q_{n,k}(1+u)\) and \(1-\Aq(y)\).  Determine whether sign-regularity,
total positivity, or continued-fraction expansions follow from the underlying
partition posets.
\end{researchproblem}

\section{Inverse generalized Gaussian monotonicity}
The integer positive-base theorem is complete, while the continuously
interpolated \(\mathcal G_q(\alpha,\beta)\) still lacks a global base-sign
diagram.

\begin{conjecture}[\statusconjectural]
For \(\alpha>\beta>0\),
\(q\mapsto\mathcal G_q(\alpha,\beta)\) is strictly increasing on \((0,1)\),
with the boundary cases \(\beta=0\) and \(\beta=\alpha\) constant.  If this is
false, the counterexample set is separated from the integer lattice by an
explicit curve and is detected by one sign change of the base derivative.
\end{conjecture}
This is intentionally phrased to invite falsification: high-precision scans
should precede any attempted proof.

\section{Inverse \texorpdfstring{\(q\)}{q}-gamma minimizer}
Beyond monotonicity of \(x_q^*\), one can ask for a complete asymptotic
expansion as \(q\to0\):
\[
 x_q^*=1+\frac{\log\lambda}{\lambda}
 +\frac{A_2(\log\lambda)}{\lambda^2}+\cdots,
 \qquad \lambda=-\log q.
\]
Near \(q=1\), implicit differentiation of \(\psi_q(x_q^*)=0\), together with
Lambert-series asymptotics, should give an ordinary power series in
\(-\log q\).  Matching these regimes may yield accurate global rational
approximants for the branch-separating minimum.

\section{Multivariate inverse systems}
Many natural identification problems contain enough observables to recover
several parameters, for example
\[
 (a;q)_\infty=y_1,
 \qquad
 \partial_a\log(a;q)_\infty=y_2.
\]
The Jacobian is built from \(S_1,S_2\), and on \(a<1\) its sign can often be
controlled.  Similar systems can recover \((q,x)\) from \((\Gamma_q(x),
\psi_q(x))\).

\begin{researchproblem}
Classify minimally sufficient collections of \(q\)-observables whose parameter
map is globally injective on natural real domains.  Develop condition numbers,
optimal experiment design, and certified multivariate inversion.
\end{researchproblem}

\section{Formalization roadmap}
A proof-assistant development can proceed in layers:
\begin{enumerate}[label=\arabic*.]
\item finite products, derivative identities, and real interval signs;
\item uniqueness of finite critical points via strict monotonicity of \(S_1\);
\item exact discriminant \eqref{eq:finite-poch-disc};
\item positive-coefficient inverse theorem for Gaussian coefficients;
\item infinite products by uniform convergence and explicit tail bounds;
\item the principal inverse through a formal inverse-function theorem;
\item asymptotic critical points using quantified estimates on \(G_m\);
\item interval certification of the \(n=5\) collision sheet labels.
\end{enumerate}
The first four layers are largely algebraic and order-theoretic and should be
well suited to Lean, Isabelle, or Coq.  The asymptotic layers require a mature
library for infinite products and uniform big-\(O\) estimates.

\chapter{Conclusion}
The inverse theory of \(q\)-analogs is not a collection of isolated numerical
root-finding tasks.  It is a coherent subject organized by four structures:
\begin{enumerate}[label=(\roman*)]
\item logarithmic derivative kernels, which control monotonicity and
conditioning;
\item discriminants and critical values, which control algebraic and analytic
branching;
\item asymptotic reversion, which converts forward \(q\)-series expansions into
inverse expansions;
\item positivity and combinatorics, which can turn a many-sheeted complex
problem into a unique global positive real branch.
\end{enumerate}

For \(q\)-Pochhammer products, every natural parameter now has a concrete
inverse framework.  The finite argument inverse has a complete real atlas and
an exact secondary-discriminant mechanism.  The infinite argument inverse has a
canonical principal branch, endpoint expansions, log-periodic large-target
behavior, and an asymptotic atlas of all remote branch points.  The base inverse
is globally monotone in the principal sign regimes and has polylogarithmic
\(q\to1\) reversion.  The order inverse reduces exactly to the principal
argument inverse plus a logarithmic lattice.

For Gaussian coefficients, nonnegative coefficients yield a particularly
strong global theorem: every target above one determines exactly one positive
base.  Bernoulli and cumulant expansions then produce systematic local inverse
series.  Continuous parameter inverses inherit monotonicity and branch geometry
from \(q\)-gamma.  The same ideas extend to \(q\)-factorials, multinomials,
\(q\)-beta, Euler exponentials, and broad basic hypergeometric families.

The most striking structural observation in this report is the appearance of Maxwell
loci: distinct simple critical points can share a critical value as the base
changes.  The exact \(n=5\) collision shows that global inverse geometry cannot
be inferred from local monotonicity or from a presumed ordering of extrema.
Secondary discriminants provide the right organizing invariant.

The next decisive advances would be a proof of the complete Bernstein
conjecture for \(1-\Aq\), a classification of finite-product Maxwell loci, and
uniform double-scaling inversions near \((a,q)=(1,1)\).  Each is accessible to a
combination of product identities, Lambert-series estimates, exact algebra, and
certified computation.

\appendix
% The article class restarts its section counter at the appendix.  Distinct
% Hyperref anchors keep appendix destinations separate from the main matter.
\renewcommand*{\theHsection}{appendix.\Alph{section}}
\renewcommand*{\theHsubsection}{\theHsection.\arabic{subsection}}
\renewcommand*{\theHsubsubsection}{\theHsubsection.\arabic{subsubsection}}
\renewcommand*{\theHtheorem}{\theHsection.\arabic{theorem}}
\renewcommand*{\theHequation}{\theHsection.\arabic{equation}}
\renewcommand*{\theHfigure}{\theHsection.\arabic{figure}}
\renewcommand*{\theHtable}{\theHsection.\arabic{table}}

\chapter{Inverse-calculus formulas}
\label{app:inverse-calculus}

\section{Derivatives of a one-variable inverse}
Let \(y=F(x)\), let \(x=G(y)\), and write
\(f_j=F^{(j)}(x)\).  At every regular point \(f_1\ne0\), repeated
implicit differentiation gives
\begin{align}
 G'&=\frac1{f_1},\label{eq:inverse-derivative-1}\\
 G''&=-\frac{f_2}{f_1^3},\label{eq:inverse-derivative-2}\\
 G'''&=\frac{3f_2^2-f_1f_3}{f_1^5},\label{eq:inverse-derivative-3}\\
 G^{(4)}&=\frac{-15f_2^3+10f_1f_2f_3-f_1^2f_4}{f_1^7},
 \label{eq:inverse-derivative-4}\\
 G^{(5)}&=
 \frac{105f_2^4-105f_1f_2^2f_3+10f_1^2f_3^2
       +15f_1^2f_2f_4-f_1^3f_5}{f_1^9}.
 \label{eq:inverse-derivative-5}
\end{align}
All derivatives on the right are evaluated at \(x=G(y)\).  These formulas
are especially effective for the principal inverse \(\Aq\), because the
alternating signs of the direct-product derivatives reduce sign questions to
explicit Bell-polynomial inequalities.

For arbitrary order, Fa\`a di Bruno's formula applied to
\(F(G(y))=y\) yields the recursion
\begin{equation}
 G^{(n)}(y)
 =-\frac1{F'(G(y))}
 \sum_{k=2}^{n}F^{(k)}(G(y))
 B_{n,k}\bigl(G'(y),G''(y),\ldots,G^{(n-k+1)}(y)\bigr),
 \label{eq:inverse-derivative-bell-recursion}
\end{equation}
where \(B_{n,k}\) is the partial exponential Bell polynomial.  Thus every
higher inverse derivative is computable without differentiating an implicit
root solver.

\section{Ordinary power-series reversion}
Suppose
\[
 f(x)=\sum_{m\ge1}f_mx^m,
 \qquad f_1\ne0,
\]
and let \(g(z)=\sum_{n\ge1}g_nz^n\) satisfy \(f(g(z))=z\).  The
Lagrange--B\"urmann formula is
\begin{equation}
 g_n=\frac1n[x^{n-1}]\left(\frac{x}{f(x)}\right)^n.
 \label{eq:lagrange-burmann-coefficients}
\end{equation}
For symbolic computation, a triangular recurrence is often faster:
\begin{equation}
 g_1=\frac1{f_1},
 \qquad
 g_n=-\frac1{f_1}[z^n]
 \sum_{m=2}^{n}f_m
 \left(\sum_{j=1}^{n-1}g_jz^j\right)^m,
 \quad n\ge2.
 \label{eq:series-reversion-triangular}
\end{equation}
Only previously known coefficients occur on the right.  Formula
\eqref{eq:series-reversion-triangular} is the all-order engine behind the
local inverse expansions in this report.

\section{Puiseux reversion at a simple fold}
Let \(F'(c)=0\), \(F''(c)\ne0\), and \(v=F(c)\).  Put
\(\eta=y-v\) and \(h=x-c\).  Expanding and reverting gives the two local
branches
\begin{align}
 h={}&\sigma\sqrt{\frac{2\eta}{F''(c)}}
 -\frac{F'''(c)}{3F''(c)^2}\eta\notag\\
 &+\sigma\frac{5F'''(c)^2-3F''(c)F^{(4)}(c)}
 {18\sqrt{2}\,F''(c)^{7/2}}\eta^{3/2}
 +O(\eta^2),
 \qquad \sigma\in\{+1,-1\},
 \label{eq:fold-puiseux-explicit}
\end{align}
with coherent choices of square roots.  A simple critical value therefore
produces quadratic ramification and a transposition in local monodromy.  At a
Maxwell point, two such folds share a target value but remain distinct source
points; at a higher caustic, one also has \(F''(c)=0\), and the Puiseux
exponent changes.

\section{Multivariate implicit inversion}
For \(F:\ComplexNumbers^m\to\ComplexNumbers^m\) and \(F(p_0)=y_0\), local inversion is possible
exactly when \(J_F(p_0)=\det DF(p_0)\ne0\).  The first derivative of the
inverse is
\[
 D(F^{-1})(y_0)=DF(p_0)^{-1}.
\]
For a scalar observable \(F(p,\lambda)=y\), with \(F_p\ne0\), variation of
an auxiliary parameter is governed by
\begin{equation}
 \frac{\partial p}{\partial\lambda}
 =-\frac{F_\lambda}{F_p}.
 \label{eq:parameter-inverse-sensitivity}
\end{equation}
This identity unifies the sign calculations for inverse base, order, and
shape parameters.  Near a branch locus the denominator warns of ill
conditioning before a numerical solver visibly fails.

\chapter{Exact secondary discriminants for finite products}
\label{app:secondary-discriminants}

For
\[
 P_n(a;q)=\prod_{j=0}^{n-1}(1-aq^j),
 \qquad
 D_n(y;q)=\Disc_a(P_n(a;q)-y),
\]
define the secondary discriminant
\[
 \mathfrak M_n(q)=\Disc_y D_n(y;q).
\]
It vanishes when two critical values coincide, when a critical point becomes
multiple, or when degree drops under a singular specialization.  Cyclotomic
and endpoint factors should therefore be separated from the factors that
encode genuine collisions of distinct finite critical values.

The exact factorizations produced by the accompanying symbolic computation
are as follows.

\begin{align}
 \mathfrak M_3(q)
 &=16q^6(q-1)^6(q^2+q+1)^3,
 \label{eq:M3-exact}\\
 \mathfrak M_4(q)
 &=16q^{42}(q-1)^{26}(q+1)^6(q^2+1)^8\notag\\
 &\quad\times
 (9q^4+31q^3+45q^2+31q+9)^3.
 \label{eq:M4-exact}
\end{align}
For \(n=5\),
\begin{align}
 \mathfrak M_5(q)
 ={}&256q^{150}(q-1)^{60}(q+1)^{14}
 (q^2+q+1)^4
 (q^4+q^3+q^2+q+1)^{15}\notag\\
 &\times
 (36q^6+97q^5+181q^4+213q^3+181q^2+97q+36)^3
 H_5(q)^2,
 \label{eq:M5-exact}
\end{align}
where
\begin{equation}
 H_5(q)=8q^{10}-5q^9-7q^8-6q^7-8q^6+35q^5
       -8q^4-6q^3-7q^2-5q+8.
 \label{eq:H5-polynomial}
\end{equation}
The reciprocal polynomial \(H_5\) has exactly one zero in \((0,1)\), namely
\[
 q_0=0.8401533783083766586629364318891329431013925836583\ldots .
\]
At this base, the critical points in the first and third zero intervals are
\begin{align*}
 c_0&=1.0714044413496100275076056970032887643067724227285\ldots,\\
 c_2&=1.5701685824623346994394512430684932917838814269399\ldots,
\end{align*}
and
\[
 P_5(c_0;q_0)=P_5(c_2;q_0)
 =-0.000295418321326036978942747465329242197437155822\ldots .
\]
Thus \(H_5=0\) is a genuine Maxwell component, not merely a degeneracy of one
critical point.

For \(n=6\), after removal of endpoint, cyclotomic, and manifestly positive
factors, one collision factor is the reciprocal degree-forty polynomial
\begin{align}
 H_6(q)={}&200q^{40}+1387q^{39}+4911q^{38}+11453q^{37}
 +18354q^{36}+17819q^{35}+69q^{34}\notag\\
 &-36849q^{33}-76446q^{32}-85091q^{31}-33909q^{30}
 +68239q^{29}+160240q^{28}\notag\\
 &+159280q^{27}+27236q^{26}-177962q^{25}-321574q^{24}
 -286466q^{23}-74232q^{22}\notag\\
 &+176918q^{21}+287794q^{20}+176918q^{19}-74232q^{18}
 -286466q^{17}-321574q^{16}\notag\\
 &-177962q^{15}+27236q^{14}+159280q^{13}+160240q^{12}
 +68239q^{11}-33909q^{10}\notag\\
 &-85091q^9-76446q^8-36849q^7+69q^6+17819q^5
 +18354q^4+11453q^3\notag\\
 &+4911q^2+1387q+200.
 \label{eq:H6-polynomial}
\end{align}
High-precision computation finds a positive zero
\[
 q=0.7773703152388489309\ldots,
\]
and continuation numerically identifies it with a collision of the critical
values from the first and third zero intervals; this sheet identification is
not interval-certified.  The complete exact factorization of
\(\mathfrak M_6\), including very large positive factors, is recorded in
\path{data/secondary_discriminants.txt}; keeping it as machine-readable
text avoids several pages of typographically uninformative coefficients.

\begin{remark}
Reciprocity of the collision factors is consistent with the reversal identity
\[
 P_n(a;q)=(-a)^nq^{n(n-1)/2}
 P_n\!\left(a^{-1}q^{1-n};q\right),
\]
together with the transformation \(q\mapsto q^{-1}\).  It does not by itself
identify which pair of critical sheets collides.  That labeling requires
continuation of the real critical points or interval-certified root isolation.
\end{remark}

\chapter{Reproducible numerical and symbolic experiments}
\label{app:reproducibility}

\section{Archive layout}
The archive contains the following source and generated artifacts.
\begin{longtable}{@{}>{\raggedright\arraybackslash}p{0.41\textwidth}
                         >{\raggedright\arraybackslash}p{0.51\textwidth}@{}}
\toprule
Path & Purpose\\
\midrule
\endhead
\path{inverse_q_analogs_report.tex}
  & Complete LaTeX source of this report.\\
\path{inverse_q_analogs_report.pdf}
  & Compiled report.\\
\path{numerical_experiments.py}
  & Documented Python program that regenerates every table, exact
    discriminant factorization, and figure used in the report.\\
\path{requirements.txt}
  & Minimal Python package requirements.\\
\path{data/finite_critical_values.csv}
  & Critical points and critical values of finite products for selected
    \((n,q)\).\\
\path{data/infinite_critical_asymptotics.csv}
  & Exact numerical critical points and heights compared with their remote
    asymptotic formulas.\\
\path{data/principal_inverse_derivative_signs.csv}
  & Derivatives of \(\Aq\) through order eight at a grid of test points.\\
\path{data/qbinomial_inverse_accuracy.csv}
  & Exact Gaussian-coefficient inverses near \(q=1\), compared with linear,
    quadratic, and cubic reversion.\\
\path{data/qgamma_minima.csv}
  & The minimizer \(x_q^*\), minimum value, and curvature of \(\Gamma_q\).\\
\path{data/secondary_discriminants.txt}
  & Exact symbolic factorizations of \(\mathfrak M_n\) for \(3\le n\le6\).\\
\path{figures/*.pdf}
  & Vector figures generated by the Python program and included in the report.\\
\bottomrule
\end{longtable}

\section{Reproduction command}
From the report directory, run
\begin{lstlisting}[style=pseudo,language=bash]
python3 numerical_experiments.py
latexmk -pdf -interaction=nonstopmode inverse_q_analogs_report.tex
\end{lstlisting}
The numerical program uses high-precision arithmetic for branch points and
root solving, exact integer polynomial arithmetic for discriminants, and
ordinary double precision only for plotting after the high-precision data have
been generated.  It fixes no random seed because no randomized algorithm is
used.

\section{Checks implemented in the program}
The script performs the following independent consistency checks.
\begin{enumerate}[label=(\roman*)]
\item Every finite critical point is bracketed between consecutive zeros and
satisfies the logarithmic-derivative equation to the working precision.
\item At the \(n=5\) collision, two separately continued critical points have
critical values agreeing to approximately fifty decimal digits.
\item Remote infinite-product critical points are compared with both the
location and height asymptotics, and the ratios approach one.
\item Gaussian inverse values are substituted into the exact finite product
formula before the asymptotic errors are recorded.
\item Derivatives of the principal product inverse are generated from a local
power-series reversion rather than finite differences.
\item The \(q\)-gamma minimizer is obtained from the zero of \(\psi_q\), and
positive \(q\)-trigamma verifies that the stationary point is a strict
minimum.
\end{enumerate}
These checks do not replace a proof of conjectural global statements, but they
make every numerical claim in the report directly auditable.

\chapter{Notation and branch conventions}
\label{app:notation}

\begin{longtable}{@{}p{0.25\textwidth}p{0.67\textwidth}@{}}
\toprule
Symbol & Meaning\\
\midrule
\endhead
\((a;q)_n\) & Finite \(q\)-Pochhammer product
\(\prod_{j=0}^{n-1}(1-aq^j)\).\\
\((a;q)_\infty\) & Infinite product for \(|q|<1\).\\
\((a;q)_\alpha\) & Complex-order continuation
\((a;q)_\infty/(aq^\alpha;q)_\infty\), with a fixed branch of \(q^\alpha\).\\
\(P_n(a;q)\) & The finite direct product \((a;q)_n\).\\
\(P(a;q)\) & The infinite direct product \((a;q)_\infty\).\\
\(\Aq(y)\) & Principal real inverse in the argument:
\((\Aq(y);q)_\infty=y\), \(\Aq(y)<1\), for \(0<q<1\), \(y>0\).\\
\(\mathcal Q_{n,k}(y)\) & Unique positive real base satisfying
\(\GaussianBinomial{n}{k}{q}=y\), for \(0<k<n\) and \(y>1\).\\
\([x]_q\) & \(q\)-number \((1-q^x)/(1-q)\), with a branch specified for
complex powers.\\
\(\Gamma_q,\psi_q\) & \(q\)-gamma and \(q\)-digamma functions.\\
\(B_q\) & \(q\)-beta function.\\
\(D_n(y;q)\) & Primary discriminant \(\Disc_a(P_n(a;q)-y)\).\\
\(\mathfrak M_n(q)\) & Secondary discriminant \(\Disc_yD_n(y;q)\).\\
Caustic & Parameter values where a critical point is degenerate or a direct
map loses ordinary fold structure.\\
Maxwell locus & Parameter values where distinct critical points have the same
critical value.\\
Principal branch & The real branch selected by a monotone interval and, in
complex continuation, by a stated base point and path.\\
\bottomrule
\end{longtable}

Unless stated otherwise, real-product results use \(0<q<1\).  Complex powers
use a locally fixed logarithm, \(q^z=\exp(z\PrincipalLogarithm q)\).  Consequently, inverses
in an exponent or order naturally repeat on a logarithmic lattice.  Square
roots in Puiseux expansions are fixed by continuation from the real side or by
an explicitly chosen local sheet.

\backmatter

\begin{thebibliography}{99}

\bibitem{ReshetnikovMonograph}
Vladimir Reshetnikov,
\emph{\(q\)-Pochhammer Symbols and \(q\)-Binomial Coefficients},
ProveIt repository, 2026.
\url{https://github.com/VladimirReshetnikov/ProveIt/blob/main/Analysis/FabiusFunction/docs/semi-formalized-research-frontiers/drafts/exponents-and-q-series/q_pochhammer_q_binomial_monograph/q_pochhammer_q_binomial_monograph.tex}.

\bibitem{DLMFqGamma}
NIST Digital Library of Mathematical Functions,
\emph{Chapter 5, Gamma Function; Section 5.18, \(q\)-Gamma and \(q\)-Beta
Functions}, edited by F. W. J. Olver et al., National Institute of Standards
and Technology.
\url{https://dlmf.nist.gov/5.18}.

\bibitem{GasperRahman}
George Gasper and Mizan Rahman,
\emph{Basic Hypergeometric Series}, second edition,
Encyclopedia of Mathematics and its Applications 96,
Cambridge University Press, 2004.

\bibitem{AndrewsAskeyRoy}
George E. Andrews, Richard Askey, and Ranjan Roy,
\emph{Special Functions},
Encyclopedia of Mathematics and its Applications 71,
Cambridge University Press, 1999.

\bibitem{McIntosh1999}
Richard J. McIntosh,
Some asymptotic formulae for \(q\)-shifted factorials,
\emph{Ramanujan Journal} \textbf{3} (1999), 205--214.

\bibitem{Moak1984}
Daniel S. Moak,
The \(q\)-analogue of Stirling's formula,
\emph{Rocky Mountain Journal of Mathematics} \textbf{14} (1984), 403--413.

\bibitem{BanerjeeWilkerson}
Debraj Banerjee and Peter Wilkerson,
Lambert series and \(q\)-functions near \(q=1\),
\emph{International Journal of Number Theory} \textbf{13} (2017), 2097--2113.

\bibitem{IsmailMuldoon}
Mourad E. H. Ismail and Martin E. Muldoon,
Inequalities and monotonicity properties for gamma and \(q\)-gamma functions,
in \emph{Approximation and Computation}, International Series of Numerical
Mathematics 119, Birkh\"auser, 1994, 309--323.

\bibitem{FlajoletSedgewick}
Philippe Flajolet and Robert Sedgewick,
\emph{Analytic Combinatorics},
Cambridge University Press, 2009.

\bibitem{Henrici}
Peter Henrici,
\emph{Applied and Computational Complex Analysis, Volume 1},
Wiley, 1974.

\bibitem{CorlessLambertW}
Robert M. Corless, Gaston H. Gonnet, David E. G. Hare,
David J. Jeffrey, and Donald E. Knuth,
On the Lambert \(W\) function,
\emph{Advances in Computational Mathematics} \textbf{5} (1996), 329--359.

\bibitem{JohanssonArb}
Fredrik Johansson,
Arb: efficient arbitrary-precision midpoint-radius interval arithmetic,
\emph{IEEE Transactions on Computers} \textbf{66} (2017), 1281--1292.

\end{thebibliography}

\end{document}
